\documentclass[10pt,a4paper]{article}

% ----------------- Paquetes -------------------%
\usepackage[T1]{fontenc}
\usepackage[utf8]{inputenc}
\usepackage[a4paper,margin=2.5cm]{geometry}
\usepackage{mathptmx}             
\usepackage{changepage} 
\usepackage{setspace}
\usepackage[spanish]{babel}
\usepackage[labelfont=bf,labelsep=space]{caption}
\usepackage{amssymb}
\usepackage{amsmath}
\usepackage{ebgaramond}
\usepackage{placeins}
\usepackage{etoolbox}
\usepackage{setspace}
\usepackage{parskip} 
\usepackage{tcolorbox}
\usepackage{needspace}
\usepackage{xparse}
\usepackage{ocgx2}
\usepackage{hyperref}
\usepackage{hypcap}


%%%%%%%%%%%%%%%%%%%%%%%%%%%%%%%%%%%%%%%%%%%%%%%%%%%%%%%%%%%%%%%%
% --- Fija primero tus valores globales "buenos" ---
\setstretch{1.2}                 % Interlineado global
\setlength{\parskip}{0.7\baselineskip}
\setlength{\parindent}{0pt}
\newlength{\GlobalParskip}
\setlength{\GlobalParskip}{\parskip}

\newcounter{eqnum}[subsection]
\renewcommand{\theeqnum}{\arabic{eqnum}}
\newcounter{alphaitem}
\newcounter{numitem}

%% ------------------- Recuadros----------------------%%
\newcommand{\puntosclave}[1]{%
  \begin{tcolorbox}[
    colframe=black,
    title={Puntos clave},
    before upper={\setlength{\parindent}{0.5cm}\setlength{\parskip}{0.5\baselineskip}}
  ]
  #1
  \end{tcolorbox}
}

\newcommand{\modificaciones}[1]{
  \begin{tcolorbox}[
    colback=white,
    colframe=red,
    title={Modificaciones a realizar},
    before upper={\setlength{\parindent}{0.5cm}\setlength{\parskip}{0.5\baselineskip}}
  ]
  #1
  \end{tcolorbox}
}

%% ------------------- Preguntas TEST----------------------%%
% Opciones: radio (single-choice)
\NewDocumentCommand{\OptRadio}{m m m}{%
  \noindent\CheckBox[radio,name=#1,value=#2,width=1.2ex,height=1.2ex]{}%
  \;\textbf{#2.}~#3\par
}

% Opciones: checkbox (multi-choice)
\NewDocumentCommand{\OptCheck}{m m}{%
  \noindent\CheckBox[name=#1,width=1.2ex,height=1.2ex,bordercolor={0 0 0}]{}%
  \;#2\par
}

% Pregunta single-choice (radio)
% Args: 1=id, 2=enunciado, 3=opciones, 4=solución+justificación, 5=imagen (o {}), 6=origen
\NewDocumentCommand{\PreguntaTestSingle}{m m m m m m}{%
  \par\medskip
  \noindent\textbf{#1.}~#2\par\smallskip

    % Imagen (si no es {})
    \IfBlankTF{#5}{}{%
      \begin{center}
        \includegraphics[width=0.75\linewidth]{#5}
      \end{center}
    }

    \begin{Form}
      #3
    \end{Form}

    \par\smallskip
    \switchocg{sol-#1}{\fbox{Mostrar / ocultar solución}}%
    \hfill{\footnotesize #6}\par\smallskip

    \begin{ocg}{Solución}{sol-#1}{0}
      \noindent #4
    \end{ocg}
  \par\medskip
}

% Pregunta multi-choice (checkboxes)
\NewDocumentCommand{\PreguntaTestMulti}{m m m m m m}{%
  \par\medskip
  \noindent\textbf{#1.}~#2\par\smallskip

    % Imagen (si no es {})
    \IfBlankTF{#5}{}{%
      \begin{center}
        \includegraphics[width=0.75\linewidth]{#5}
      \end{center}
    }
    \begin{Form}
      #3
    \end{Form}

    \par\smallskip
    \switchocg{sol-#1}{\fbox{Mostrar / ocultar solución}}%
    \hfill{\footnotesize #6}\par\smallskip

    \begin{ocg}{Solución}{sol-#1}{0}
      \noindent #4
    \end{ocg}
  \par\medskip
}

%% ------------------- Listas----------------------%%
    % --- Lista alfabética ---
    \newenvironment{la}
      {\begin{list}{\alph{alphaitem})}{%
          \usecounter{alphaitem}%
          \setlength{\leftmargin}{1cm}%
          \setlength{\parskip}{3pt}% 
          \setlength{\itemsep}{0pt}% ← espacio entre ítems
          \setlength{\parsep}{0pt}%  ← párrafos dentro de un ítem
      }}
      {\end{list}}
      
    % --- Lista numérica ---
    \newenvironment{lnum}
      {\begin{list}{\arabic{numitem}.}{%
          \usecounter{numitem}%
          \setlength{\leftmargin}{1cm}%
          \setlength{\parskip}{3pt}% 
          \setlength{\itemsep}{0pt}% ← espacio entre ítems
          \setlength{\parsep}{0pt}%  ← párrafos dentro de un ítem
      }}
      {\end{list}}
      
% ------------------- Imágenes----------------------%
    \graphicspath{{Img/}}% Rutas por defecto 
    % --- Imagen adaptativa ---
    % Registros de dimensión definidos una sola vez
    \newdimen\imgcap
    \newdimen\imgmin
    \newdimen\imgmax
    \newdimen\imgremain
    \newdimen\imgh
    \newdimen\imgneed
    
    \newcommand{\imagenfit}[3]{%
      \addvspace{10pt}% separación antes
      \par\begingroup
        % Parámetros de composición (asignaciones, no \newdimen)
        \imgcap=1\baselineskip            % alto estimado de título+fuente
        \imgmin=0.15\textheight
        \imgmax=0.15\textheight
        \imgremain=\pagegoal
        \ifdim\pagegoal=\maxdimen \imgremain=0pt\fi
        \advance\imgremain by -\pagetotal
        \advance\imgremain by -\imgcap
        % Decisión de altura destino
        \ifdim\imgremain<\imgmin
          \newpage
          \imgh=0.20\textheight
        \else
          \imgh=\imgremain
          \ifdim\imgh>\imgmax \imgh=\imgmax \fi
        \fi
        % Reservar espacio para evitar cortes del bloque completo
        \imgneed=\imgh \advance\imgneed by \imgcap
        \Needspace{\imgneed}
        % Cuerpo
        \noindent\begin{minipage}{\textwidth}
          \centering
          \captionof{figure}{\textemdash\ #2}\par\vspace{0.3em}% título arriba
          \includegraphics[width=\linewidth,height=\imgh,keepaspectratio]{\detokenize{#1}}\par
          \vspace{0.3em}\small\textit{Fuente: }#3% fuente abajo
        \end{minipage}%
      \endgroup
      \par\addvspace{10pt}%
    }

%------------ Bloque de RECUADRO ESPECIAL -------------%
    \newenvironment{specialblock}[1]{%
      \begin{quote} % crea sangría a izquierda y derecha
      \small % texto más pequeño
      \noindent\textbf{#1}\\[-0.3em] % título en negrita
      \rule{\linewidth}{0.4pt}\\[0.6em]}{
      \end{quote}}
      
%-------------------------- Portada -----------------------%
\newcommand{\oldtitle}[1]{\def\OldTitle{#1}}
\newcommand{\changerationale}[1]{\def\ChangeWhy{#1}}

\newcommand{\changebox}{%
\begin{tcolorbox}[title={Historial de cambios del tema}]
\textbf{Título anterior:} \OldTitle\par
\textbf{Motivo del cambio:} \ChangeWhy
\end{tcolorbox}
}

%-------------------------- Author Font -----------------------%
    \newcommand{\authorfont}[1]{{\fontfamily{EBGaramond-TLF}\selectfont #1}}

%-------------------------- Anexos -----------------------%
    \makeatletter
    \newcounter{anexo}
    \renewcommand{\theanexo}{\Roman{anexo}}
    \newcommand{\anexo}[1]{%
      \clearpage
      \phantomsection
      \refstepcounter{anexo}%
      \noindent\textbf{Anexo \theanexo.- }#1\par\medskip
      \addcontentsline{stc}{section}{Anexo \theanexo.- #1}%
    }
    \makeatother

    %-------------------------- Lista de figuras (personal) -----------------------%
\makeatletter
\newcommand{\MyListOfFigures}{%
  \clearpage
  \phantomsection
  {\centering\bfseries\Large Índice de figuras\par}%
  \medskip
  % Entrada en nuestro índice de contenido (stc)
  \addcontentsline{stc}{section}{Índice de figuras}%
  % Recuperación del listado (sin cabecera propia)
  \begingroup
    \parindent=0pt\parskip=2pt
    \@starttoc{lof}%
  \endgroup
}
\makeatother

%-------------------------- Bibliografía -----------------------%
\makeatletter
% Contador para las referencias
\newcounter{myref}
% Cabecera de la bibliografía, entra en el índice 'stc'
\newcommand{\MyBibliography}{%
  \clearpage
  \phantomsection
  {\centering\bfseries\Large Bibliografía y referencias\par}%
  \medskip
  \addcontentsline{stc}{section}{Bibliografía y referencias}%
  \setcounter{myref}{0}%
}
\newcommand{\MyBibItem}[4]{%
  \refstepcounter{myref}%
  \noindent[\themyref]\ #1.\ \emph{#2}. #3. #4\par
}
\makeatother

%--------- Establecer cuatro niveles de índice --------%
    \setcounter{tocdepth}{4}   
    \setcounter{secnumdepth}{4}% numera hasta \paragraph

%%%%%%%%%%%%%%%%%%%%%%%%%%%%%%%%

\makeatletter

    %--------- Interlineado Global --------%
    \edef\GlobalBaseLineStretch{\baselinestretch}
    
    %--------- Espaciado de seccinones --------%
    \renewcommand\section{\@startsection{section}
      {1}{\z@}
      {2.5ex \@plus 1ex \@minus .2ex} % espacio antes
      {1.5ex \@plus .2ex}             % espacio después
      {\normalfont\Large\bfseries}    % formato del título
    }
    \renewcommand\subsection{\@startsection{subsection}{2}{\z@}
      {3ex \@plus 0.8ex \@minus .2ex}
      {1.5ex \@plus .2ex}
      {\normalfont\large\bfseries}
    }
    \renewcommand\subsubsection{\@startsection{subsubsection}{3}{\z@}
      {3ex \@plus 0.5ex \@minus .2ex}
      {1ex \@plus .2ex}
      {\normalfont\normalsize\bfseries}
    }
    \renewcommand\paragraph{\@startsection{paragraph}{4}{\z@}%
    {-2ex \@plus -0.5ex \@minus -.2ex}% espacio antes
    {0.8ex \@plus .2ex}% espacio después
    {\normalfont\normalsize\itshape}}% ← cursiva en lugar de negrita

    %--------- Ecuaciones --------%   
    % Variables globales para almacenar títulos y notas
    \gdef\leq@current@section{}%
    \gdef\leq@current@subsection{}%
    \gdef\leq@last@printed@section{}%
    \gdef\leq@last@printed@subsection{}%
    
    % Comando para añadir nota (invisible en el cuerpo)
    \newcommand{\eqnote}[1]{%
      \addcontentsline{leq}{leqnote}{#1}%
    }
    
    % Comando para ecuaciones
    \newcommand{\eqblock}[2]{%
      % Verificar si necesitamos imprimir el título de sección
      \ifx\leq@current@section\leq@last@printed@section\else
        \addcontentsline{leq}{leqsection}{\leq@current@section}%
        \global\let\leq@last@printed@section\leq@current@section
        \gdef\leq@last@printed@subsection{}% Reset subsección al cambiar sección
      \fi
      % Verificar si necesitamos imprimir el título de subsección
      \ifx\leq@current@subsection\leq@last@printed@subsection\else
        \ifx\leq@current@subsection\@empty\else
          \addcontentsline{leq}{leqsubsection}{\leq@current@subsection}%
          \global\let\leq@last@printed@subsection\leq@current@subsection
        \fi
      \fi
      % Ecuación
      \refstepcounter{eqnum}%
      \begin{equation*}
        #1\tag{E.\theeqnum}
      \end{equation*}%
      \addcontentsline{leq}{leqt}{[E.\theeqnum] -- #2}%
      \addcontentsline{leq}{leqm}{#1}%
    }
    
    %%% ----- Índice de ecuaciones ----------%%%
    % Tamaño para []
    \newcommand{\leqIndexFont}{\normalsize}
    \newcommand{\leqTitleFont}{\tiny}
    \newcommand{\leqNoteFont}{\tiny}
    % Cabecera de sección 
    \newcommand*\l@leqsection[2]{%
      \addvspace{25pt}% Mayor separación antes de nueva sección
      \noindent\leqIndexFont
      \makebox[\linewidth]{\rule{.38\linewidth}{0.4pt}\ \ \textbf{  \MakeUppercase{#1}}\ \ \rule{.38\linewidth}{0.4pt}}%
      \par\vspace{-10pt}
    }
    % Cabecera de subsección 
    \newcommand*\l@leqsubsection[2]{%
      \par\addvspace{15pt}%
      \noindent\leqIndexFont #1
      \par\vspace{-7pt}
    }
    % Título de ecuación 
    \newcommand*\l@leqt[2]{%
    \par\addvspace{10pt}%
      \par\noindent\leqTitleFont #1\par
    }
    % Miniatura de ecuación
    \newcommand*\l@leqm[2]{%
      \vspace{0.3\baselineskip}%
      \par\scriptsize\noindent\hspace*{1.5em}\(#1\)\par%
    }
    % Nota de ecuación 
    \newcommand*\l@leqnote[2]{%
      \vspace{0.2\baselineskip}%
      \par\noindent\leqNoteFont\hspace*{1.5em}\textbf{Nota.--} #1\par\vspace{0.5\baselineskip}%
    }
    % Generación del índice
    \newcommand{\mylistofequations}{%
      \clearpage
      \phantomsection
      % Título visual de la lista (ajústalo a tu gusto):
      {\centering\bfseries\Large Índice de ecuaciones\par}
      % Entrada en nuestro índice 'stc':
      \addcontentsline{stc}{section}{Índice de ecuaciones}%
      % Llamada a tu lista real (usa el nombre exacto de tu macro):
      \@starttoc{leq}%
    }
    
    %--------- Captura de títulos de sección --------%
    \let\leq@original@section\section
    \let\leq@original@subsection\subsection
    % Flag para saber si estamos en una sección asterisco
    \newif\ifLeqStarSection
    \LeqStarSectionfalse
    
    % Redefinir \section CON CONTROL DE ASTERISCO
    \renewcommand{\section}{%
      \@ifstar{\leq@section@star}{\@dblarg\leq@section@nostar}%
    }
    \def\leq@section@star#1{%
      \global\LeqStarSectiontrue
      \gdef\leq@current@section{#1}%
      \gdef\leq@current@subsection{}%
      \leq@original@section*{#1}%
      \global\LeqStarSectionfalse
    }
    \def\leq@section@nostar[#1]#2{%
      \global\LeqStarSectionfalse
      \gdef\leq@current@section{#2}%
      \gdef\leq@current@subsection{}%
      \leq@original@section[#1]{#2}%
    }
    % Redefinir \subsection
    \renewcommand{\subsection}{%
      \@ifstar{\leq@subsection@star}{\@dblarg\leq@subsection@nostar}%
    }
    \def\leq@subsection@star#1{%
      \global\LeqStarSectiontrue
      \gdef\leq@current@subsection{#1}%
      \leq@original@subsection*{#1}%
      \global\LeqStarSectionfalse
    }
    \def\leq@subsection@nostar[#1]#2{%
      \global\LeqStarSectionfalse
      \gdef\leq@current@subsection{#2}%
      \leq@original@subsection[#1]{#2}%
    }

    % ----- Restauración de valores Globales de estilo ----------%
    % Flags
    \newif\ifInFootnote
    \InFootnotefalse
    \pretocmd{\@footnotetext}{\InFootnotetrue}{}{}
    \apptocmd{\@footnotetext}{\InFootnotefalse}{}{}
    % Profundidad de listas (soporta anidadas)
    \newcount\MyListDepth \MyListDepth=0
    \AtBeginEnvironment{la}{\advance\MyListDepth by 1}
    \AfterEndEnvironment{la}{\advance\MyListDepth by -1}
    \AtBeginEnvironment{lnum}{\advance\MyListDepth by 1}
    \AfterEndEnvironment{lnum}{\advance\MyListDepth by -1}
    \AtBeginEnvironment{itemize}{\advance\MyListDepth by 1}
    \AfterEndEnvironment{itemize}{\advance\MyListDepth by -1}
    \AtBeginEnvironment{enumerate}{\advance\MyListDepth by 1}
    \AfterEndEnvironment{enumerate}{\advance\MyListDepth by -1}
    % Restauración diferida: hazlo al INICIO del próximo párrafo real
    \newif\if@pendingRestore
    \@pendingRestorefalse
    \newcommand{\RequestRestore}{\global\@pendingRestoretrue}
    \everypar{%
      \if@pendingRestore
        \global\@pendingRestorefalse
        \ifInFootnote\else
          \ifnum\MyListDepth=0
            \setlength{\parskip}{\GlobalParskip}%
            \setstretch{\GlobalBaseLineStretch}%
          \fi
        \fi
      \fi
    }
    % Al final de bloques:
    \AfterEndEnvironment{equation}{\RequestRestore}
    \AfterEndEnvironment{equation*}{\RequestRestore}
    \AfterEndEnvironment{la}{\RequestRestore}
    \AfterEndEnvironment{lnum}{\RequestRestore}
    % Al final de macros:
    \apptocmd{\eqblock}{\RequestRestore}{}{}
    \apptocmd{\imagenfit}{\RequestRestore}{}{}
    
\makeatother

% ===================== ToC propio con 4 niveles (único bloque) =====================
\makeatletter

% --- Escritores: añadimos una línea al .stc en cada cabecera numerada ---
% Guardamos las versiones actuales para llamar al comportamiento normal
\let\MyTOC@orig@section\section
\let\MyTOC@orig@subsection\subsection
\let\MyTOC@orig@subsubsection\subsubsection
\let\MyTOC@orig@paragraph\paragraph

% SECTION
\renewcommand{\section}{%
  \@ifstar{\MyTOC@section@star}{\@dblarg\MyTOC@section@nostar}%
}
\def\MyTOC@section@star#1{%
  \MyTOC@orig@section*{#1}% sin entrada al .stc
}
\def\MyTOC@section@nostar[#1]#2{%
  \MyTOC@orig@section[{#1}]{#2}% comportamiento normal (escribe al .toc)
  % Escribimos también nuestra línea al .stc (texto con \numberline)
  \addcontentsline{stc}{section}{\protect\numberline{\thesection}#2}%
}

% SUBSECTION
\renewcommand{\subsection}{%
  \@ifstar{\MyTOC@subsection@star}{\@dblarg\MyTOC@subsection@nostar}%
}
\def\MyTOC@subsection@star#1{%
  \MyTOC@orig@subsection*{#1}%
}
\def\MyTOC@subsection@nostar[#1]#2{%
  \MyTOC@orig@subsection[{#1}]{#2}%
  \addcontentsline{stc}{subsection}{\protect\numberline{\thesubsection}#2}%
}

% SUBSUBSECTION
\renewcommand{\subsubsection}{%
  \@ifstar{\MyTOC@subsubsection@star}{\@dblarg\MyTOC@subsubsection@nostar}%
}
\def\MyTOC@subsubsection@star#1{%
  \MyTOC@orig@subsubsection*{#1}%
}
\def\MyTOC@subsubsection@nostar[#1]#2{%
  \MyTOC@orig@subsubsection[{#1}]{#2}%
  \addcontentsline{stc}{subsubsection}{\protect\numberline{\thesubsubsection}#2}%
}

% PARAGRAPH
\renewcommand{\paragraph}{%
  \@ifstar{\MyTOC@paragraph@star}{\@dblarg\MyTOC@paragraph@nostar}%
}
\def\MyTOC@paragraph@star#1{%
  \MyTOC@orig@paragraph*{#1}%
}
\def\MyTOC@paragraph@nostar[#1]#2{%
  \MyTOC@orig@paragraph[{#1}]{#2}%
  % Si no numeras los \paragraph, cambia \theparagraph por texto plano sin \numberline
  \addcontentsline{stc}{paragraph}{\protect\numberline{\theparagraph}#2}%
}

% --- Impresor del índice propio (lee el .stc) ---
\newcommand{\MyTOC}{%
  \par\bigskip
  {\centering\bfseries\Large Índice de contenido\par}%
  \medskip
  \begingroup
    \parindent=0pt\parskip=2pt
    \@starttoc{stc}%
  \endgroup
}

\makeatother
% ===================== fin ToC propio =====================


%%%%%%%%%%%%%%


% --- Datos del tema ---
\title{Unión Europea (II)\\
\large Las finanzas de la Unión Europea y el presupuesto comunitario. El marco financiero plurianual actual}
\author{Marcelo Ortega}
\date{Marzo 2026}
\newcommand{\TemaCode}{3B37}

\oldtitle{
    B.38. Las finanzas de la Unión Europea y el presupuesto comunitario. El marco financiero plurianual actual
}
\changerationale{
    Sin cambios.
}

\usepackage{soul}\sethlcolor{yellow}% resaltado amarillo: \hl{texto} (Panel TCEE, ⌃H)
\begin{document}


\maketitle
\changebox

\newpage

% --- Página de resumen y modificaciones ---
\puntosclave{ %% Escribir puntos clave Apuntes CECO
     A la hora de realizar la exposición oral del tema se recomienda dedicar 5 minutos a la introducción y a la conclusión, 17-20 minutos a la PAC dando especial relevancia a la última reforma y a su configuración actual y 5-8 minutos a la PPC.

     Si bien se recoge de manera exhaustiva el proceso de reforma de las políticas agraria y pesquera, se recomienda al opositor que, en el momento de la exposición, recorra los antecedentes históricos de manera ejecutiva y se centre en la situación actual de las mismas.

     Las ideas clave que deberían transmitirse son: ¿Por qué son necesarias estas políticas? ¿Cómo se articularon inicialmente y cuáles fueron sus resultados? ¿Qué factores han motivado la necesidad de refonnarlas? ¿valoración de las mismas? ¿Qué rumbo seguirán estas políticas y qué otras problemáticas condicionarán su desarrollo futuro?
    }
\vspace{1cm}

\modificaciones{
    \textbf{NOTA (04/10/2026):} Revisar: https://www.politico.eu/article/european-parliament-wants-eu-long-term-budget-to-top-e2-trillion/
}

\newpage
\MyTOC
\newpage

\mylistofequations
\clearpage

\MyListOfFigures
\clearpage


\pagenumbering{arabic}
\setcounter{page}{1}


\section{Introducción}



\subsection{Contextualización}


\subsubsection{Relevancia}




\subsubsection{Problemática}



\subsection{Estructura de la exposición}





\clearpage
\section{Política Agraria Común}
\puntosclave{
    Consultar: \href{https://www.elconfidencial.com/mundo/reflexiones-europeas/2026-07-03/politica-agraria-comun-europa-1hms_4381996/}{\textit{La Política Agrícola Común en su dimensión histórica}. El Confidencial (2026)}; \href{https://www.elconfidencial.com/mundo/2026-07-03/ucrania-moldavia-agricultura-ue-integracion-1hms_4381734/}{\textit{¿Pánico agrícola u oportunidad en la UE? La próxima batalla de Ucrania será el grano}. El Confidencial (2026)}; \href{https://www.elconfidencial.com/economia/2024-03-24/funciona-la-pac-los-agricultores-no-tienen-razon_3854235/}{\textit{¿Funciona la PAC? Los agricultores no tienen razón (salvo en alguna cosa)}. El Confidencial (2025)}; \href{https://www.elconfidencial.com/espana/tribuna/2026-04-24/recorte-pac-gobierno-espana-europa-1hms_4343699/}{\textit{El futuro de la PAC en manos del peor Gobierno}. El Confidencial (2026)}; \href{https://www.elconfidencial.com/espana/2024-02-06/que-es-pac-ley-protesta-huelgas-agricultores_3825271/}{\textit{¿Qué es la PAC y por qué es importante en la huelga de agricultores?} El Confidencial (2026)}

    La PAC responde a la necesidad de garantizar la correcta implementación del Mercado Interior en el ámbito agrícola. 
    
    Política dinámica, se ha adaptado a los sucesivos desafios de la agricultura europea.
    
    La última reforma busca garantizar un futuro sostenible al sector. prestar un apoyo más específico a las explotaciones más pequeñas y ampliar la flexibilidad de los Estados miembros a la hora de adaptar las medidas a las condiciones locales. 
    
    Sus principales elementos son, un cambio de orientación estratégica. una mayor ambición medioambiental y la inclusión, por primera vez, de una dimensión social.
}

La PAC se creó en 1962 con los objetivos de mejorar la productividad agrícola para que los consumidores dispusiesen de un suministro estable de alimentos a precios asequibles y garantizar a los agricultores un nivel de vida razonable. En su concepción inicial, la PAC se articulaba a través de un sistema de ayudas vinculadas a la prodttcción que favorecieron un sistema de producción intensivo generador de grandes excedentes, daños medioambientales, evolución explosiva del gasto agrícola, así como discriminaciones entre regiones, productos y explotaciones y distorsiones en el comercio internacional. Para corregir estos desequilibrios, la PAC se hasometido a sucesivas reformas y ha registrado una profunda transformación hasta alcanzar su configuración actual basada en un sistema de ayudas multifuncionales de apoyo a las rentas y redes de seguridad para los productores, en el que se ha reforzado la orientación al mercado, el desarrollo rural y los criterios medioambientales. Las sucesivas reformas también han permitido hacer frente a los desafíos a los que se ha ido enfrentando el sector agrícola europeo como perturbaciones l!n el mercado mundial, la utilización más sostenible de los recursos naturales, la mitigación del cambio climático, garantizar la prosperidad de las zonas rurales y la seguridad alimentaria. De manera específica, la última reforma de la PAC aspira a garantizar un futuro sostenible a los agricultores y ganaderos europeos, prestar un apoyo más específico a las explotaciones más pequeñas y ampliar la flexibilidad de los Estados miembros a la hora de adaptar las medidas a las condiciones locales.

A nivel nacional, el sector agrícola se ha visto tradicionalmente sometido a una minuciosa regulación debido a su peso económico y social, su debilidad estructural y su sensibilidad coyuntural. La agricultura es uno de los sectores económicos más sensibles debido a su vulnerabilidad a tres tipos de fenómenos: naturales (riesgos meteorológicos, aleatoriedad de las cosechas, dependencia del clima, carácter perecedero de las producciones), económicos (variabilidad de los precios, gran competencia en el sector, baja elasticidad renta de los alimentos) y sociales (número de explotaciones familiares, envejecimiento de los activos, inmovilidad de la fuerza de trabajo, formación limitada).

La necesidad de una política agrícola se ha visto posteriormente reforzada al tomar consciencia de los efectos externos negativos derivados de ciertas prácticas agrícolas y de los bienes públicos que los agricultores proveen al desarrollar otras funciones más allá de la tradicional provisión de alimentos como son la gestión sostenible de los recursos naturales, la conservación del paisaje y el mantenimiento de una economía rural viva.

A nivel comunitario, la implementación del Mercado Interior y la libre circulación de productos agrícolas manteniendo al mismo tiempo la intervención pública, hacía necesario suprimir los mecanismos de intervención nacionales y transferirlos a nivel supranacional para poder gestionar los eventuales efectos negativos de la integración sobre el campo y evitar que las políticas nacionales introdujesen distorsiones en el Mercado Interior.

\subsection{¿Qué es y por qué está? Naturaleza y objetivos}
El fundamento jurídico de la PAC se encuentra recogido en los artículos 38 a 44 del Tratado de Funcionamiento de la Unión Europea (TFUE). El artículo 4.2, letra d) del TFUE. establece una competencia compartida entre la UE y los Estados miembros en el ámbito agrícola.

La configuración inicial de la PAC en el Tratado de Roma respondía a un contexto posbélico, caracterizado por la escasez de alimentos en Europa tras la Segunda guerra mundial. En ese contexto, se fijaron unos objetivos específicos de carácter económico y social (Art. 39 TFUE) con los que se pretendía proteger los intereses de los productores (aumentar la productividad, asegurar un nivel de vida equitativo) y de los consumidores (asegurar los aprovisionamientos y precios razonables). Varias disposiciones del TFUE incorporan otros objetivos aplicables al conjunto de las políticas de la UE, y que, en consecuencia, son también objetivos de la PAC, relacionados con el empleo y la cohesión económica, social y territorial, la protección del medio ambiente y de los consumidores, el bienestar animal y la salud pública.

En 1958 tuvo lugar la Conferencia lntergubernamental de Stressa para definir las líneas directrices de la PAC. El resultado de esta Conferencia fue el Plan Mansholt 1, aprobado en 1962, que recogía las proposiciones de la Comisión para la progresiva edificación de la PAC. Entre sus proposiciones, la Comisión señaló tres puntos esenciales sobre los que debía basarse la política agraria:
• Organización de los mercados agrarios por tipos de productos.
• Mejora de las estructuras agrarias.
• Ordenación de los intercambios con terceros países.

De acuerdo a las propuestas de la Comisión, la PAC quedó articulada en tomo a·dos políticas, la Política de precios y de mercados y la Política de estructuras agrarias. De relevancia desigual, la primera ha sido tradicionalmente la más importante en cuanto a actuaciones y recursos absorbidos, mientras que la segunda ha tenido escasa entidad debido a la insuficiencia de los recursos asignados y la exigencia de cofinanciación por parte de los Estados Miembros.

\subsubsection{Principios}
También en la Conferencia lntergubername.ntal de Stressa se establecieron los principios que debían regir el funcionamiento de la PAC que fueron adoptados en 1962:

• Principio de Unidad de Mercado: garantizaba la libre circulación de los productos agrarios dentro de la Comunidad, excluyendo cualquier discriminación entre los productores o consumidores de cualquiera de los Estados miembros, la fijación de precios idénticos para todo el territorio comunitario y establecía una estricta regulación del mercado interior.\footnote{%
    Se garantizaba mediante la supresión i1-1terior de barreras arancelarias y no arancelarias. las restricciones cuantitativas. la supresión de ayudas nacionales, un sistema de tramitación administrativa común de las exportaciones e importaciones de productos agrarios. así como la coordinación y armonización de las medidas sanitarias.
}
• Principio de Preferencia Comunitaria: suponía la protección de la producción interna respecto de la competencia exterior mediante un complejo sistema de defensa en frontera con el objetivo de favorecer la producción, el comercio y el consumo de los productos comunitarios. Permitía el mantenimiento de precios más elevados para la producción agrícola comunitaria que los vigentes en el mercado mundial mediante el establecimiento de un tratamiento exterior común en la gestión de las importaciones y de las exportaciones.
• Principio de Solidaridad Financiera: implicaba la financiación de los gastos derivados de la PAC a través del presupuesto comunitario, con la salvedad de I.a Política de estructuras agrarias que sería cofinanciada por los Estados miembros. A su vez, los ingresos generados por la PAC tenían la consideración de recursos propios de la Comunidad.


\subsection{Desarrollo histórico: Proceso de consolidación}


\subsubsection{Configuración inicial: Política de precios y Política de Estructuras Agrarias}
La financiación de la PAC corría a cargo del Fondo Europeo de Orientación y Garantía A grícola (FEOGA), creado en 1962 y separado en dos secciones en 1 964. • La sección Garantía, con mucho la más significativa, tenía por objeto la financiación de los gastos derivados de la aplicación de la Política de precios y de mercados. Esta sección financiaba íntegramente las medidas de intervención en los mercados. • La sección Orientación financiaba las operaciones de la Política de estructuras agrarias. Esta sección se basaba en el Principio de cofinanciación.


\paragraph{Política de precios y de mercados}
La Política de precios y de mercados se articuló a través de las Organizaciones Comunes de Mercados (OCMs), disposiciones relativas a la producción, comercialización y tratamiento de los productos agrícolas en la UE. Las OCMs regulaban la política común de precios y los mecanismos de intervención. Los precios de los productos se fijaban artificialmente al principio de cada campaña de comercialización, distinguiéndose tres tipos de precios:
Precios Guía: En cada OCM recibían un nombre distinto: precio indicativo, que se aplicaba a aquellos productos cuyas cotizaciones son precisas y se pueden almacenar (cereales, azúcar, leche, aceite de oliva y semillas de colza y de girasol); de orientación, para los productos cuyas cotizaciones no son muy precisas y en su mayor parte son perecederos (bovinos pesados, vinos y ciertos pescados); de base, para la carne de porcino y ovino y objetivo para el tabaco, algodón, lino ... 

Garantizados o precios suelo : En la OCM de cereales, los organismos de intervención tenían la obligación, con las limitaciones de cantidades y calidades establecidas, de adquirir los productos ofertados por los agricultores a estos precios que recibían el nombre de precios de intervención. En la OCM de frutas y hortalizas, cuando se alcanzaban estos precios denominados de retirada, se concedían ayudas para retirar la producción del mercado hasta que los precios se recuperasen. Una actuación similar se producía en la OCM del vino, en la que los precios garantizados se denominaban desencadenantes. 

De entrada o precios techo: Igual que en los casos anteriores, según el producto recibían un nombre diferente: umbral (cereales, azúcar, mantequilla, aceite de oliva), esclusa (porcino, aves y huevos) o de referencia (frutas y hortalizas).

Para fijar precios idénticos para los productos agrícolas en todos los países miembros, se creó en 1962 sólo a efectos de los precios agrarios la Unidad de Cuenta Agrícola (UCA), definida por una paridad de oro igual que la del dólar, de tal fonna que los precios se establecían en UCAs y después se transfonnaban a cada moneda nacional aplicando el tipo de cambio vigente. La utilización de la UCA no generó graves inconvenientes hasta la devaluación del franco francés y la revaluación del marco alemán en 1969 que provocó un incremento de los precios agrícolas en Francia y una reducción de los mismos en Alemania, además de un incremento de las ayudas percibidas por los agricultores franceses y una disminución de las ayudas obtenidas por los agricultores alemanas. Como solución, se pennitió a Francia y Alemania mantener en el sector agrario sus antiguos tipos de cambio frente a la UCA. A partir de ese momento se distinguieron dos tipos de cambio, el tipo de conversión o tipo verde, aplicable exclusivamente a la agricultura, y el tipo central, aplicable al resto de la economía. Los tipos verdes se fijaban para cada campaña coincidiendo con el establecimiento de los precios guía y pennanecían al margen de las fluctuaciones de los tipos centrales. Para evitar que bajo este sistema se produjesen movimientos especulativos se crearon los montantes compensatorios monetarios (MCM), gravámenes y subvenciones que se aplicaban al comercio intracomunitario y de terceros países de ciertos productos agroalimentarios, principalmente a los intercambios de aquellos productos en los que las compras de intervención eran relevantes en la regulación. Actuaban como un impuesto sobre las exportaciones para los países cuya moneda se había devaluado y como una subvención para las monedas que se habían revaluado. Al adherirse un nuevo Estado miembro a la Comunidad, durante el periodo transitorio y para adecuar sus precios a los comunitarios, se aplicaban a los intercambios entre la Comunidad y el nuevo Estado miembro los montantes compensatorios de adhesión (MCA). Los MCA adoptaban la fonna de un gravamen o una subvención cuya cuantía era equivalente a la diferencia entre los precios de los miembros recién llegados y los comunitarios. La existencia de los MCM, afiadía una gran complejidad técnica, administrativa y financiera, tanto para los servicios comunitarios y de los Estados como para los importadores y exportadores agrarios. Este sistema agromonetario intentó desmantelarse sin éxito en varias ocasiones. La implementación del Mercado Único y la abolición de las fronteras intracomunitarias obligaron a eliminar este sistema y desde enero de 1993 los tipos verdes se ajustaban a los centrales y se eliminaron los MCM.

Compras por los organismos de intervención: Se distinguían dos tipos de OCMs, con garantía total de precios y con garantía parcial de precios. En las primeras, la intervención era obligatoria,a requerimiento del productor, cualquiera que fuese la situación del mercado y siempre dentro de los meses establecidos. La compra se efectuaba al precio guía. En las segundas, las compras sólo estaban previstas en situaciones de descenso de los precios de mercado. Éstas situaciones podían ser definidas de manera precisa en la reglamentación correspondiente o decididas por la Comisión o el Consejo. La compra podía limitarse a ciertas zonas de la Comunidad, a ciertas cantidades y a los precios de compra que eventualmente se determinasen.

Todas las OCMs establecían mecanismos de intervención entre los que destacaban:
• Compras por los organismos de intervención de los productos que les eran ofrecidos por los productores en las condiciones establecidas en la reglamentación del sector. 

• Ayudas al almacenamiento privado: mecanismo de retirada del producto del mercado de manera temporal. Esta medida se orientaba a conseguir la reducción de excesos de oferta coyunturales, sin necesidad de incurrir en los elevados costes de las compras de intervención.

• Ayudas a la producción: subvenciones a los productores que propiciaban el descenso de los precios en origen y que, en consecuencia, incentivaban la demanda de los productos tanto para consumo como para su utilización por la industria.

• Ayuda a la eliminación de excedentes: estas medidas buscaban la desnaturalización de un producto excedentario de manera que se destinase a usos industriales. El ejemplo más relevante era la destilación de vino para la obtención de alcohol. La Comunidad completaba esta política con un régimen común de intercambios con terceros países que propiciaba la separación entre el mercado comunitario y el mercado mundial consiguiendo un nivel de precios en el interior de la Comunidad más elevado que en el mercado mundial.

• Medidas a la importación: Uno de los elementos era el Arancel exterior común (AEC). No obstante, era insuficiente para garantizar que los productos de países terceros no entrasen en la Comunidad a precios inferiores al precio mínimo de importación. Para evitarlo, se aplicaban exacciones reguladoras y tasas compensatorias. o Exacciones reguladoras (prélévements). Mecanismo de protección comunitaria por excelencia y principal garante del Principio de preferencia comunitaria. Se trataban de gravámenes variables, normalmente específicos. equivalentes a la diferencia entre el precio del producto en el mercado comunitario y en el mercado mundiát. o Tasas compensatorias o montantes suplementarios. Se calculaban según el precio de oferta de cada expedición comercial y no según el precio del mercado mundial.

• Medidas a la exportación: o Restituciones a la exportación. Su objetivo era compensar la diferencia entre el precio del producto en el mercado comunitario y en el mercado mundial, generalment�� más bajo, para facilitar la exportación de los productos excedentarios. o Las exacciones a la exportación y la suspensión de las exportaciones. Ambas medidas podían ser adoptadas cuando los precios del mercado mundial eran superiores a los existentes en el mercado comunitario, de manera que existía incentivo a exportar y dejar desabastecido el mercado comunitario.

Para algunos cultivos como las materias grasas y proteaginosas o las semillas oleaginosas, en los que la Comunidad era deficitaria, se utilizaba el sistema de deficiency payments. Para asegurar los abastecim_ientos, la producción comunitaria se incentivaba mediante subvenciones a la producción que compensasen la diferencia ente los precios comunitarios y los del mercado mundial. Los productos se intercambiaban en el mercado comunitario a los precios del mercado mundial, y el agricultor recibía un ingreso igual a la diferencia entre su coste de producción y el precio del mercado.

\paragraph{Política de estructuras agrarias}
La Política de estructuras agrarias se definió en el · Plan Mansholt !, aunque fue con el segundo Plan Mansholt ( 1 969), cuando se sentaron las bases de su funcionamiento. Algunos de los objetivos de la PAC como asegurar un nivel de vida equitativo a los agricultores, unos precios razonables a los consumidores y un equilibrio entre la oferta y demanda de los productos agrícolas no podía garantizarse con la única aplicación de una Política de precios y mercados por tres razones:
• Las enormes diferencias de productividad de la economía europea. La Política de precios y mercados aseguraba a las explotaciones bien estructuradas y económicamente rentables unas rentas excesivamente elevadas, mientras que las explotaciones marginales obtenían unos ingresos insu��cientes.
• Un régimen de precios en este contexto estructural no podía asegurar unos precios razonables a los consumidores porque si los precios debían posibilitar la subsistencia de explotaciones marginales, estos debían ser elevados en comparación con los que regirían si sólo funcionaran explotaciones modernas.
• Un régimen de precios aplicado sobre esta realidad estructural no podía conducir a un equilibrio entre oferta y demanda porque los altos precios inducidos por esta política eran un incentivo a la sobreproducción. 

Algunas de las principales medidas contempladas por esta política fueron: 
• De mejora de la eficacia de las estructuras agrarias: o Ayudas a la inversión en explotaciones agrícolas. Se concedían a explotaciones en las que la renta del trabajo por unidad de trabajo/hombre era inferior a una renta de referencia establecida por el Estado miembro. Se perseguía la mejora cualitativa y la reconversión de la producción en función de las necesidades del mercado, la diversificación de actividades en las explotaciones, la reducción de los costes de producción y la mejora de las condiciones de vida y de trabajo. o Ayudas a la instalación de jóvenes agricultores. En forma de primas o bonificación de intereses que se concedían para subvencionar los gastos de primera instalación y los planes de mejora de jóvenes agricultores.  
• Medidas específicas en beneficio de la agricultura de montaña y zonas desfavorecidas para mejorar las condiciones de transformación y comercialización de los productos.

\paragraph{Resultados}
La PAC se desarrolló entre 1 962 y 1 969. En 1 965 ya estaba regulada el 85\% de la producción final agraria y en 1 968 la PAC entraba en su fase definitiva con la aplicación de precios comunes. Se logró la reestructuración del sector favoreciéndose el crecimiento, mejoras en la gestión y la adquisición de conocimientos técnicos por los agricultores. Entre 1969 y 1984 la PAC se consolidó, lográndose los objetivos fijados en el Tratado de Roma, aunque al mismo tiempo se detectaron los siguientes efectos perversos que dieron lugar a una pérdida de legitimidad interna y externa del modelo agrícola europeo: 
• Un aumento espectacular de la productividad. Se alcanzaron niveles de producción superiores al nivel de absorción del mercado debido a la falta de adecuación de la oferta a la demanda. Parte de los grandes excedentes generados se destinaban al interior de la Comunidad que para su mantenimiento y eliminación consumían una parte sustancial del FEOGA-Garantía: Otra parte se destinaba al exterior en forma de exportaciones subsidiadas mediante restituciones a la exportación que permitían compensar la diferencia entre el precio interior y el del mercado mundial. Esta circunstancia generó tensiones con terceros países, especialmente con los países en desarrollo. 
• Producción intensiva con efectos perjudiciales para el medioambiente, seriamente cuestionada por la opinión pública. 
• Dinámica explosiva del gasto agrícola. La PAC llegó a representar más del 70\% del presupuesto comunitario. La práctica totalidad del gasto recaía sobre los consumidores y evitaba que se abordasen otras políticas comunitarias. • lnequidad en la distribución de las ayudas entre :o Sectores y Estados miembros: Los productos atlánticos (cereales, leche, vacuno) estaban más favorecidos que los mediterráneos (vid, olivo). Los cereales y la ganadería concentraban el 70\% de los gastos agrícolas totales. o Explotaciones: al determinarse las subvenciones en función de los rendimientos, las grandes explotaciones resultaron más beneficiadas. 

En 1 969 la Comisión publicó el "Memorandum sobre la reforma de la PAC", también conocido como Plan Mansholt 11, con el que se inició un debate que concluiría con la aprobación en 1 972 de una serie de medidas socioestructurales encaminadas a:
• Fomentar el crecimiento y modernización de las explotaciones más eficientes.
• Reducir los costes que provocaría la reestructuración, fundamentalmente relacionados con la expulsión de la mano de obra sobrante y la eliminación de las explotaciones no viables.
• Reducir los costes de producción para conseguir abaratar la política de sostenimiento de precios.

Sus resultados fueron decepcionantes como consecuencia de la crisis económica de la década de 1970, el incumplimiento de los plazos establecidos para el desarrollo normativo por parte de los Estados miembros y la falta de concreción en las regiones más atrasadas debido a la necesidad de cofinanciación de las medidas.

\subsubsection{Reforma: Libro Verde sobre las perspectivas de la PAC (1985)}

Identificadas las debilidades del modelo inicial, la Comisión presentó el Libro Verde sobre las Perspectivas de la: PAC en 1985, que planteaba una serie de medidas para una eventual reforma.
a) Objetivos perseguidos
• Disminuir la producción excedentaria mediante una política de mercados que permitiera ajustar la oferta a la demanda.
• Frenar el crecimiento de los gastos presupuestarios.
• Potenciar el desarrollo de las áreas rurales para preservar la función de conservación del entorno de la agricultura y promover la protección del medio ambiente.
b) Principales medidas adoptadas.
• Medidas limitativas de la producción: Tasa de corresponsabilidad, reducción del período de intervención, arranque de viñas y cuotas en el sector lácteo.\footnote{%
    Tasas de corresponsabilidad: Se fijaron cantidades máximas garantizadas en el sector de los cereales. En caso de ser superadas, los productores debían pagar una tasa de corresponsabilidad equivalente al 3\% del precio de intervención. La superación de las cantidades garantizadas también implicaba una reducción del precio de intervención para la campaña siguiente en un 3\%.

Arranque de viñas: Los Estados miembros debían obligatoriamente otorgar una subvención por hectárea a los viticultores que aceptasen de manera voluntaria arrancar con carácter definitivo sus viñas.
}
• En 1988 se introdujeron estabilizadores presupuestarios, mecanismo relativamente automático de aplicación en las OCMs más relevantes que consistía en la reducción de precios y garantías cuando la producción o la intervención superaban determinados límites. En las perspectivas financieras del periodo 1988-1992 se introdujo la llamada directriz agrícola que limitaba el crecimiento del gasto agrícola al 74% del crecimiento del PIB comunitario.

• Medidas socioestructurales orientadas a disminuir la oferta y mejorar las estructuras agrarias y el medio ambiente. o Retirada de tierras arables. Se otorgó una subvención a los agricultores por retirar sus tierras o parte de las mismas del cultivo para reducir la producción. o Extensificacion de cultivos. Esta medida pretendía limitar la producción de excedentes reduciendo la intensividad de la producción a cambio de una subvención. o Incentivo a la prejubilación. Programa voluntario para los Estados consistente en otorgar una ayuda a los agricultores de más de 55 años que decidiesen cesar definitivamente en la actividad agrícola.

\textcolor{red}{La reforma de los fondos estructurales de 1 988 incorporó la Política de estructuras agrarias en la Política de cohesión económica y social. Dicha reforma concentró las ayudas de fondos ivos. Las zonas rurales se integraron en el objetivo I y en el obJ<- o 5b) 4 . ..1s explotaciones agrarias susceptibles de ayudas estructurales para moderrnzarse y 111,::J-,1 dr su productividad se incluyeron en el objetivo 5a), centrado en acelerar la adaptación de las estructuras agrarias. Dentro de las iniciativas comunitarias se desarrolló el programa Leader para av,ufar a las asociaciones rurales o gobiernos locales en dichas zonas a potenciar la actividad ,ii versificación.}

\subsubsection{Reforma de MacSharry: Comunicación sobre Evolución y Futura de la PAC (1991)}
A pesar de .la reforma de 1985 que permitió controlar el crecimiento de los gastos de 1, PAC y
redujo los excedentes en algunos productos, persistían numerosos problemas que fueron
analizados por la Comisión en la Comunicación sobre Evolución y Futuro de [:; P/\C ( i 991).
• Nuevos incrementos en los excedentes de vacuno, leche, vino y tabaco.
• Las medidas de apoyo a la retirada de tierras sólo habían afectado al 2% de las tianias de
cultivo de cereales. El resto de medidas de acompañamiento prácticamente no se
aplicab!ln.
• Constatación de que el 80% de las ayudas del FEOGA se destinaban al "' 1
-; de las
explotaciones.
• Estancamiento de la renta agraria a pesar de que la pot-J,
un 35% en el período 1975-89.
• Constatación del aumento de los ¡
• Tensiones con terceros países '
,1cida en
��xcede .61 1-. ,1...s en los
mercados mundiales. Este inci. ..: ofe. ,_ .... ·ficticia"' causó una caída en los  precios agrícolas mundiales que agravó las tensiones presupuestarias en la Comun-i:d��d.La Comunicación de la Comisión se materializó en un:'l nropuesta de refonna, l::i refünna MacSharry, aprobada en 199:? grícolas siguientes. 

a) Objetivos perseguidos
• Lograr una agricultura europea competitiva en el mercado internacional.
• Desarrollar una agricultura respetuosa con el medioambiente.
b) Principales medidas adoptadas
• Reducir sustancialmente los precios garantizados de forma gradual.
• Introducir ayudas directas por hectárea o por cabeza de ganado, teniendo en cuenta los rendimientos históricos, como compensación a la caída de los ingresos de los agricultores derivada de la reducción de precios.
• Para el cobro de estas ayudas debían realizarse retiradas obligatorias de tierras arables de forma rotatoria y fomentar la extensificación de la agricultura.
• En las perspectivas financieras acordadas para el período 1993- 1999, se mantuvo la directriz agrícola.
• Medidas socioestructurales concentradas en mejorar y proteger el medioambiente. Se introdujo un verdadero programa agroambiental en el que las ayudas iban destinadas a los agricultores que se comprometieran a adoptar o mantener prácticas que mejoraran o preservaran el medio ambiente (limitar el uso de productos químicos, introducir o continuar con métodos orgánicos, mantener las tierras arables abandonadas y los bosques.)

El apoyo a las zonas rur��les y a la mejora de las explotaciones agrarias en la Política de Cohesión y Social prosiguió, de tal forma que el nuevo Reglamento de los Fondos estructurales para el período 1993-99, mantuvo los objetivos de la primera reforma ampliándose a las zonas de pesca.

Tras la Ronda Uruguay del Acuerdo General sobre Aranceles Aduaneros y Comercio (GA TI), el sistema tradicional de protección frente a terceros países desaparece. Al establecerse un sistema de protección basado exclusivamente en aranceles de carácter fijo, los precios de entrada dejaron de tener valor y los prélévements, las tasas compensatorias y los montantes suplementarios perdieron su carácter variable y se integraron en la estructura arancelaria de la Comunidad. Las subvenciones a la exportación fueron disminuyendo progresivamente y se establecieron sistemas más eficaces de control de la producción agraria para evitar los excedentes. Respecto al sistema de ayudas a los agricultores, estas se vieron reducidas y modificadas con la introducción de ayudas directas por hectárea y por cabeza de ganado, que se incluyeron en la "Caja azul". 

El desarrollo de las negociaciones comerciales de la Ronda Uruguay fue uno de los factores que motivó la refoma de 1992. Los compromisos asumidos en el Acuerdo de agricultura en materia  de acceso a mercados, ayuda interna y subvenciones a las exportaciones modificaron la configuración inicial de la PAC.
✓ En acceso a mercados. los países se comprometieron a convertir en aranceles las barreras comerciales no arancelarias y establecer niveles máximos para los mismos. Igualmente, se comprometieron a reducir los niveles arancelarios máximos en una media del 36% con un mínimo por línea del 15% en el período 1995-2000.
✓ Respecto a la ayuda interna, se identificaron tres tipos de ayudas en fu nción de su nivel de distorsión sobre el comercio internacional. Las ayudas más distorsionantes (ayudas a la industria de transformación, siempre que beneficiasen a los productores del producto de base; la ayuda al sostenimiento de precios de mercado; las ayudas directas vinculadas a la producción y las subvenciones a los inputs y ayudas a los costes de comercialización) se incluyeron en la denominada "Caja ámbar". Los países debían calcular la Medida Global de Ayuda (MGA) correspondiente al período 1986-88 y reducir su importe en un 20% en tramos anuales iguales. En contraposición a la "Caja ámbar", las ayudas que no distorsionaban el comercio o lo hacían mínimamente se englobaron en la denominada "Caja verde" (programas gubernamentales (fonnación, investigación, e infraes tructuras, ayuda alimentaria interna y existencias públicas para seguridad) y pagos directos a los productores producidos por catástrofes y seguros, desarrollo rural, políticas de medio ambiente y ayudas directas a las rentas no vinculadas a la producción.) En la "Caja azul" se incluyeron de los pagos directos a los agricultores y ganaderos establecidos en el marco de programas destinados a limitar la producción, que de manera transitoria se excluyeron del cálculo de la MGA.
✓ Por último, las exportaciones subvencionadas debían disminuirse un 21 %, tomando como referencia los períodos 1986- 1989 o 1990-92 según productos, y el gasto presupuestario en subvencio!:)eS a la exportación debía reducirse en un 36%.

La reforma MacSharry coincidió con la Cumbre de la Tierra de Río en la que ��e reclamó un desarrollo más sostenible con el objetivo de alcanzar un equilibrio entre el desarrollo de la actividad agrícola y el respeto al medioambiente en el que el agricultor desempeña un papel relevante a través de su función de gestión del espacio y protección de los recursos naturales. Se introduce a partir de este momento el concepto de multifuncionalidad de la actividad agraria, al reconocer que los agricultores no sólo producen alimentos, sino que, además, desempeñan labores de protección medio-ambiental y de ordenación del territorio. 

De la reforma MacSharry se pueden destacar aspectos positivos y negativos. Entre los primeros, la reducción de los excedentes (cereales y vacuno), la disminución en términos relativos del coste de la PAC y el aumento de la renta agrícola media un 4,5\% en el periodo 1992- 1996.\footnote{%
    Crisis de las vacas locas en el ganado inglés o los problemas de las dioxinas en la carne de pollo belga.
} Hasta entonces la PAC había estado financiada básicamente por los consumidores a través de las subvenciones ocultas en los precios. pero con la refonna. serían los contribuyentes los que sostendrían ese coste. Con la introducción de ayudas directas explícitas también se logró mayor transparencia en las actuaciones de la política agraria al conocerse los destinatarios y los montos recibidos. Entre los segundos. se reprodujeron algunos de los viejos problemas en cuanto a distribución territorial, sectorial y empresarial de las ayudas. Al haberse calculado las ayudas en función de los rendimientos históricos, las zonas más productivas consiguieron las primas más elevadas y se consagró el trato favorable otorgado a los países del norte (con mayor productividad) en detrimento de los del sur (con más bajos rendimientos). Por último, se mantuvieron las grandes diferencias entre los precios comunitarios y los mundiales. al mismo tiempo que la protección en frontera seguía siendo elevada. 

\subsubsection{Reforma \textit{abierta}: }
La reforma de 1992 sentó las bases para solucionar parte de los desequilibrios de la PAC y en los años siguientes, se realizaron continuos ajustes englobados en la llamada «reforma abierta» articulada en tres fases: Agenda 2000, la de 2003 y el llamado «chequeo médico» (2008-2009).

\paragraph{Agenda 2000: }
La agricultura europea se enfrentaba de nuevo a retos internos y externos.

• Los excedentes se incrementaron de nuevo en los últimos años de la década de los noventa y los estudios realizados por la Comisión en 1999 preveían un empeoramiento de la  situación en el horizonte 2005/2006. A esta situación se añadía el estudio de los mercados agrarios de los países de Europa Central y Oriental, candidatos a la adhesión a la VE, que preveía incrementos en los excedentes.
• Era necesario mejorar el cumplimiento del compromiso adquirido en el Tratado de Ámsterdam de integrar la protección del medioambiente en toda la legislación comunitaria.
• Las crisis alimentarias6 obligaron a los legisladores a adoptar un nuevo enfoque que reforzara la seguridad alimentaria. A esta circunstancia se sumaba el reciente interés de los consumidores por productos alimenticios seguros y de mayor calidad y las reclamaciones de un mayor bienestar para los animales.
• Necesidad de una mayor transparencia, descentralización y simplicidad.
• La ampliación de los países del Este aumentaría la superficie cultivada en un 55% y la población agraria activa en un 1 3%.
• En el año 2000 se iniciaron fonnalmente unas nuevas negociaciones agrícolas en el seno de la OMC que conllevarían la reducción de la protección comunitaria, la reducción o eliminación de las subvenciones a la expo11ación y la reducción de las ayudas internas, incluida la negociación sobre el mantenimiento de las ayudas directas contabilizadas dentro de la Caja azul.

La Comisión presentó en 1997 el documento Agenda 2000: por una Unión más fuerte y más amplia y la reforma propuesta fue aprobada en 1999. 

a) Objetivos perseguidos
Lograr.una agricultura competitiva a nivel interior y exterior, que garantizase al mismo un nivel de vida equitativo y unas rentas estables para los agricultores.
• Diseñar una agricultura duradera y de calidad, que respetase el medio ambiente y ofreciese productos de calidad que respondiesen a las expectativas de los consumidores.
• Implantar una agricultura que preservase y potenciase el mundo rural, en línea con el concepto de multifuncionalidad de la agricultura.
• Conseguir una política más cercana a la población de forma que los costes de esa política fuesen legítimos por los servicios que los agricultores prestan a la sociedad.
• Formular una política simplificada que estableciese con claridad cuáles debían ser los costes comunes y los que debían reservarse a los Estados miembros.
b) Principales medidas adoptadas
• Se realizó un nuevo ajuste de los precios internos con respecto a los precios mundiales. Las pérdidas se compensaron mediante un aumento parcial de las ayudas directas por hectárea o cabeza de.ganado.
• Se introdujeron los principios de Ecocondicionalidad y Modulación de las medidas.o Ecocondicionalidad: Los Estados miembros establecerían ciertas medidas medioambientales a las que deberían ajustarse los agricultores y ganaderos para recibir las ayudas incurriendo en penalizaciones en caso de no hacerlo.o Modulación: Una vez delimitados unos criterios unitarios para no distorsionar la competencia, los Estados podrían modular las ayudas según criterios nacionales (nivel global de renta, el empleo de la explotación o el volumen de la ayuda monetaria recibida). Se perseguía acercar más la PAC a las realidades concretas de cada Estado. Las ayudas  que no se destinaban a las explotaciones podían utilizarse para financiar medidas de desarrollo rural. El impacto de esta medida fue limitado debido a su carácter voluntario.
• Se suprimieron o limitaron los regímenes de intervención sustituyéndolos por ayudas para el almacenamiento temporal.

• En línea con las conclusiones de la Conferencia Europea sobre Desarrollo Rural "La Europa Rural. Perspectivas de Futuro" (Cork 1996), se reforzaron las medidas estructurales, en particular las medidas agromedioambientales, en el seno de una nueva Política de desarrollo rural, heredera de la Política de estructuras agrarias. Con el - objetivo de mantener los territorios rurales vivos y desarrollar su economía potenciando sus recursos específicos se estableció una política de desarrollo rural cuya misión fuese completar las políticas de mercado y sufragar los gastos de mejora del medio ambiente y de ordenación del espacio. A partir de este momento la PAC se articula a través de la Política de precios y mercados denominada Primer Pilar y la nueva Política de desarrollo rural, denominada Segundo Pilar.

Se fijaron como objetivos específicos de esta nueva Política: o Crear un sector agrícola y forestal o Incrementar la competitividad de las áreas rurales. o Mantener el entorno y preservar la herencia rural de Europa. Los principios fundamentales sobre los que se basó la nueva Política rural fueron: o Multifuncionalidad de la agricultura. o Enfoque multisectorial e integrado. o Flexibilización de las ayudas, basándose en el principio de subsidiaridad.

o Transparencia en la .elaboración y gestión de los programas de desarrollo rural. En estos programas se diseñaba una estrategia plurianual que pretendía responder tanto a las necesidades específicas de los Estados miembros (o las regiones) como a las prioridades de la Política europea de desarrollo rural, a partir de un «menú» de medidas europeas. Se previeron medidas de acompañamiento y de modernización y diversificación de las explotaciones agrícolas. Dentro de las primeras, se añadieron las indemnizaciones compensatorias en las zonas desfavorecidas o sujetas a limitaciones medioambientales. La novedad respecto a la normativa anterior es que las medidas agroambientales pasan a ser de obligado cumplim iento para los Estados miembros. Entre las segundas, se introdujeron ayudas para inversiones en las explotaciones, instalación de jóvenes agricultores, formación profesional, mejora de la transformación y comercialización de los productos agrícolas, silvicultura y fomento de la adaptación y desarrollo de las zonas rurales. 

Al estar integrados los programas de desarrollo rural en la Política de Cohesión Económica y Social, la reforma de esta poi ítica en la Cumbre de Berlín de 1999 introdujo ciertas novedades: o Las zonas rurales (antes objetivo 5b) se mantuvieron dentro del nuevo objetivo 2. o Dentro de las zonas rurales se distingue las que están dentro de las denominadas zonas menos favorecidas (objetivo 1) y el resto. Las primeras seguirían recib iendo recursos del Feoga-Orientación, mientras que las segundas recibirían el apoyo para el desarrollo rural del Feoga-Garantía. o Se aprobó el programa Leader + dotado con 2020 millones € para el periodo 2000-2006. A diferencia de los anteriores programas, podían acceder a él todas las zonas rurales.o El antiguo objetivo 5a) (mejora de las estructuras agrarias) se incorpora en la política agrícola general y será financiado por el Feoga-Garantía.
• Se mantuvo la estabilidad presupuestaria mediante la inclusión de la directriz agrícola en las perspectivas financieras 2000-2006. En relación con los países de la ampliación se aprobó el programa Sapard para el periodo 2000-2006 para la adhesión y la preadhesión.
• Se intensificaron las nonnas de calidad y seguridad alimentaria. Se desarrollaron los principios y requisitos generales que debían cumplir las leyes alimentarias, creando también la Agencia europea para la seguridad alimentaria (EFSA). La creciente importancia de los factores distintos del precio en las decisiones de compra de los consumidores llevó a la UE a establecer normas que difundiesen de manera adecuada la calidad de los productos: normas de comercialización, sellos de calidad Denominación de origen protegida (DOP), Indicación geográfica protegida (IGP) y un logotipo especial para los productos de la agric ultura ecológica.\footnote{%
    Normas de comercialización: Se aplican a la mayoría de los productos agrícolas y definen categorías de productos.nonnas mínimas para estos y algunas condiciones de etiquetado. Proporcionan información al consumidor sobre el origen o la variedad, por ej emplo y le penniten comparar precios entre productos de calidad similar. Una DOP califica a un producto alimenticio que se ha producido en su totalidad dentro de una zona geográfica definida con técnicas e ingredientes reconocidos de la región y que tiene un vínculo con su origen geográfico. Una IGP califica a un producto vinculado por su calidad y reputación a una región en la que se desarrolla al menos una etapa de la producción. El logotipo especial para los productos de la agricultura ecológica garantiza que se han cumplido las nomias europeas de producción ecológica.
}

\paragraph{Reforma de 2003: }
a) Objetivos perseguidos. La reforma estuvo motivada por la necesidad de dar respuesta urgente
a dos cuestiones: • La futura ampliación a I O nuevos países. • Las negociaciones agrícolas en el marco de la Ronda Doha de la OMC.

b) Principales medidas (Anexo 3)
• Cambios en el sistema de ayudas directas. o Principio de desacoplamiento o desvinculación. Se disociaron las ayudas respecto de los volúmenes producidos a partir de 2005. Estas ayudas disociadas adoptaron la forma de un pago único por explotación cuyo importe quedaba determinado en función de la media de los importes que cada explotación recibió en el período histórico de referencia 2000- 2002. Los agricultores dispondrían de una gran flexibilidad para producir cualquier tipo de producto, sin que los pagos estuviesen ligados a un cultivo o a - una determinada especialización ganadera en concreto. Ante el riesgo de un abandono masivo de determinados cultivos se introdujeron matizaciones y excepciones parciales de tal forma que puede hablarse de una desvinculación sólo parcial, puesto que determinados tipos de ayudas siguieron ligadas a cultivos y producciones concretas. o Condicionalidad de la recepción de las ayudas al cumplimiento de varios requisitos medioambientales, de sanidad y bienestar animal, salubridad alimentaria y seguridad en el trabajo. En caso de incumplimiento las ayudas se reducirían en proporción al daf'lo causado. o Aplicación obligatoria de la modulación.

\begin{specialblock}{Principio de desacoplamiento}
    La introducción del Principio de desacoplamiento de las ayudas, permitió el desplazamiento del grueso del apoyo interno prestado por la PAC desde la Caja azul a la Caja verde. El modelo de aplicación de estos pagos ofrecía varias posibilidades según el grado de desacoplamiento aplicado y según las ayudas se concediesen sobre la base individual de lo percibido por el agricultor en el período de referencia 2000-2002 o se optase por un sistema regional izado:
    ✓ Fórmula de base (histórica): cada agricultor recibiría derechos correspondientes a las ayudas obtenidas durante el período de referencia y al número de hectáreas cultivadas durante dicho período que dieron derecho a ayudas directas. 
    ✓ Fórmula regional (a tanto alzado): los importes de referencia se calculaban a escala regional y correspondían a la suma de las ayudas recibidas por los agricultores en la región correspondiente durante el período de referencia. Los importes de referencia regionales se dividían por el número de hectáreas admisibles declaradas por los agricultores en el año deintroducción del Régimen de Pago Único (RPU) a fin de determinar el valor de un solo derecho en dicha región. Cada agricultor recibía un número de derechos (a tanto alzado) igual al número de hectáreas admisibles declaradas en el año de introducción del RPU.
    ✓ Modelos mixtos: en casos justificados, los Estados miembros podían aplicar distintos métodos de cálculo en las diversas regiones de su. territorio. También podían calcular las ayudas del RPU utilizando una fórmula en parte histórica y en parte a tanto alzado. Las fórmulas híbridas de este tipo podían ir variando a lo largo del período entre la aplicación inicial del RPU y la plena aplicación, dando lugar a fórmulas híbridas dinámicas y estáticas. Podía recurrirse a las fórmulas híbridas dinámicas para pasar del modelo de base (histórico) al regional (a tanto alzado).
    
    Los países miembros podían introducir ayudas adicionales (nacionales o regionales) limitadas a un máximo del' 10% de las ayudas recibidas por cada explotación para favorecer detenninados cultivos por razones medioambientales, promoción de la calidad o la comercialización.
    
    En los primeros años posteriores a la adhesión, los nuevos Estados miembros podían optar por un Régimen de pago único por superficie consistente en el pago de importes uniformes por hectárea, hasta alcanzar el límite máximo nacional establecido en los Acuerdos de adhesión. 
    
    Respecto a la aplicación obligatoria de la modulación, las explotaciones que percibiesen más de 5.000€ anuales en ayudas directas verían reducido su importe en un 3\% en 2005, aumentando al 4\% en 2006 y al 5\% a partir de 2007 y el dinero ahorrado se destinaría al desarrollo rural. De este dinero, el 80\% sería gastado en el Estado miembro donde tuviese lugar la modulación y el resto se dividiría entre los Estados miembros sobre la base de criterios de cohesión social y características estructurales de su agricultura. Los Estados miembros podrían establecer retenciones adicionales para constituir reservas nacionales de derechos de pago o bien para disponer de fondos para llevar a cabo pagos extra en determinadas zonas geográficas, principal mente por moti vos ambientales. Las regiones ul traperiféricas ( entre I as que se encuentran las Islas Canarias) no estarían sujetas ni a modulación ni a desconexión.  

    Respecto a la Política de desarrollo rural se introdujeron nuevas medidas que podían ser cofinanciadas en el marco de los programas de desarrollo rural: ✓ Incentivos a los agricultores para mejorar la calidad de sus productos y de los procesos de producción. ✓ Ayudas temporales para ayudar a los agricultores a adaptarse a las nuevas normas comunitarias, en ternas medioambientales, de seguridad en el trabajo, de bienestar del ganado, siempre que aún no hayan sido incorporadas a la legislación nacional. ✓ Establecimiento de un nuevo sistema de asesoramiento a los agricultores, con costes cubiertos principalmente por el sector público. ✓ Ayudas a los agricultores que acepten compromisos de al menos cinco aflos para usar técnicas tendentes a mejorar el bienestar de los animales de granja. ✓ Ayudas a la inversión destinadas a jóvenes agricultores. ✓ Incorporación de una medida que abre la posibilidad de otorgar pagos a superficies de uso agrario y forestal incluidas en la Red Natura. La Red Natura integra los espacios ecológicamente valiosos designados para su protección en la Directiva Aves (79/409/CEE) y en la Directiva Hábitats (92/43/CEE). En estas zonas la actividad agraria puede verse sometida a restricciones técnicas y limitaciones dando lugar a una merma en la rentabilidad que la medida de desarrollo rural mencionada trataría de compensar.
\end{specialblock}

• Fortalecimiento de la Política de desarrollo rural que registró un importante aumento de recursos financieros y una ampliación de sus objetivos e instrumentos. El nuevo Reglamento de Desarrollo Rural de 2005 promovió la preparación de programas nacionales y regionales de desarrollo rural organizados en bloques de medidas dentro de cuatro grandes líneas de actuación:
o Eje 1 : Aumento de la competitividad del sector agrícola forestal.
o Eje 2: Mejora del medioambiente y del entorno rural.
o Eje 3: Calidad de vida en las zonas rurales y diversificación de la economía rural.
o Eje 4: Aplicación del enfoque Leader, estrategias de desarrollo rural y de cooperación interterritorial basadas en grupos de iniciativas locales.

En líneas generales la reforma incrementó la orientación al mercado favoreciendo la flexibilidad en la elección de cultivos u orientaciones productivas por parte del agricultor. Permitió también compaginar el mantenimiento del objetivo de apoyo a las rentas con una menor distorsión del comercio internacional .. Sin embargo, los pagos directos por explotación siguieron siendo; aunque en menor medida, un instrumento de apoyo a la producción y al estar basados en las ayudas históricamente recibidas por cada explotación, se mantuvieron los niveles de mayor apoyo a las grandes explotaciones y a la agricultura de los países del Norte de Europa. 

En el Consejo de Bruselas de 2002 se fijó el límite de los recursos comunitarios destinados a fi nanciar la PAC para el periodo 2007-20 1 3 de acuerdo al Principio de disciplina financiera. Se congeló el presupuesto de la Política de precios y mercados y se impusieron límites anuales obligatorios.

Desde el I de enero de 2007 el FEOGA es reemplazado por el FEADER (Fondo Europeo Agrícola de Desarrollo Rural), destinado a cofinanciar los programas de desarrollo rural y el FEAGA (Fondo Europeo de Garantía Agrícola) que tiene a su cargo la financiación de las intervenciones en los mercados, los pagos directos y las subvenciones a la exportación.

También en 2007 se creó una única Organización Común de los Mercados, que codificó los mecanismos de reglamentación de las 21 OCM existentes. El objetivo fue introducir más sencillez y transparencia en la legislación de los 2 1 productos, que hasta 2007 se regían por OCMs individuales. proporcionando un marco jurídico único para el mercado interior, los intercambios con terceros países y las normas de competencia. Al derogar los reglamentos de base anteriores, se consiguió una simplificación técnica y normas más homogéneas que reemplazaron los planteamientos sectoriales por otros de carácter horizontal.

\paragraph{Chequeo médico (2009): }
En 2007 la Comisión realizó una valoración de la implementación de la refonna de 200 3 y en noviembre de 2008 se aprobó "el chequeo médico" de 2009, que permitió revisar un amplio abanico de medidas aplicadas tras la refonna de 2003.
a) Objetivos perseguidos : modernizar,- simplificar y racionalizar la PAC, eliminando las restricciones impuestas a los agricultores para permitirles reaccionar mejor ante las tendencias del mercado y hacer frente a las nuevas dificultades de la agricultura europea como el cambio climático, la necesidad de gestionar mejor el agua o la conservación de la biodiversidad.

b) Principales medidas
• Reforzar la disociación total de las ayudas mediante la eliminación progresiva de los últimos pagos asociados a la producción.
• Reorientar parcialmente los fondos de la Política de precios y mercados a favor de la Política de desarrollo rural, aumentado el índice de modulación de las ayudas directas.
• Flexibilizar las normas de intervención pública y de control de la oferta. Determinados mecanismos de intervención se vieron suprimidos y otros endurecieron sus principios de actuación.
• Flexibilizar la condicionalidad suprimiendo las normas no pertinentes e introduciendo nuevos requisitos como una mejor gestión del agua.
• Posibilitar la concesión de recursos adicionales a regiones desfavorecidas, actividades agrarias de tipo vulnerable y para apoyar medidas de gestión del riesgo como los regímenes de seguro en caso de catástrofes y los fondos de inversión para enfermedades de los animales.


\subsubsection{Propuesta: PAC en el horizonte 2020 (2010)}
La Comisión presentó en 201 O la Comunicación "La PAC en el horizonte 2020: Responder a los retos futuros en el ámbito territorial. de los recursos naturales y alimentarios•·. adoptada formalmente en 2013 y que supuso, por primera vez, a una revisión completa de la PAC.

a) Objetivos perseguidos: contribuir a la consecución de las prioridades de la Comunicación de la Comisión "Europa 2020: Una estrat��gia para un crecimiento inteligente. sostenible e integrador", así como hacer frente a los siguientes retos:

• Económicos: responder a las crecientes preocupaciones relativas a la seguridad alimentaria en la UE y en el mundo, mantener y mejorar su competitividad en un mundo cada vez más globalizado, caract!!rizado por un aumento de la volatilidad de los precios, a la vez que mantener la producción agrícola y hacer frente al deterioro de la posición de los agricultores en la cadena alimentaria.

• Medioambientales: promover la gestión sostenible de los recursos naturales, hacer frente a la presión cada vez mayor sobre las condiciones de producción agrícola causada por el cambio climático, así como a la necesidad de que los agricultores reduzcan su contribución a las emisiones de gases de efecto invernadero, desempeñen un papel activo en la mitigación del cambio climático y proporcionen energías renovabl��s.

• Territoriales: aprovechar la diversidad de las estructuras y de los sistemas de producción y mantener su papel social, territorial y de cohesión especialmente a través de la creaciónde empleo.

b) Principales medidas.
• Se suprimen todas las medidas de control de la oferta.

• Se introducen medidas excepcionales para evitar perturbaciones del mercado (fluctuaciones de los precios u otros acontecimientos), medidas de apoyo en caso de enfermedades animales y pérdida de confianza de los consumidores debido a la existencia de riesgos para la salud pública o la sanidad de los animales o las plantas, y medidas relativas a las prácticas concertadas durante los períodos de desequilibrios graves en los mercados. De esta forma se consolidan las herramientas de la OCM única como «redes de seguridad».

• Se crea una reserva para crisis que afecten a la producción o la distribución. Se constituye aplicando cada año una reducción de los pagos directos en el marco del mecanismo de disciplina financiera. En caso de que la reserva no se utilice, su importe se restituye cada año a los agricultores. La reserva de crisis puede utilizarse para financiar las medidas excepcionales.

• Las disposiciones relativas a las organizaciones de productores (OP), a las asociaciones de organizaciones de productores (AOP) y a las organizaciones interprofesionales (OIP) se amplían a todos los sectores para reforzar el poder de negociación de los agricultores.

• Se mantienen las restituciones a la exportación únicamente para algunos productos y cuando las condiciones del mercado interior se correspondan con las determinadas por las medidas excepcionales.

• El sistema de ayudas a los agricultores se configura como un sistema de pagos multifuncional, por niveles o estratos, en el que cada componente de la ayuda está vinculado a objetivos específicos y los periodos de referencia históricos no desempeñan ningún papel. Este sistema presenta componentes obligatorios y voluntarios: 
    Componentes obligatorios: Régimen de pago básico o pago simplificado por superficie, armonizado en función de criterios económicos o administrativos, a escala nacional o regional, y sometido a un proceso de "convergencia interna" por el que las asignaciones basadas en referencias históricas deben pasar a sistemas de pagos por hectárea. Ecologización. Para compensar los costes asociados al suministro de bienes públicos medioambientales no remunerados por el mercado, cada explotación recibirá un pago adicional por hectárea por respetar determinadas prácticas agrícolas beneficiosas para el clima y el medio ambiente.\footnote{%
        Diversificación de cultivos. mantenimiento de los pastos pennancntes existentes. mantenimiento de una «superficie de interés ecológico» de por lo menos el 5\% de las tien·as de labor de la explotación para aquellas explotaciones con más de 15 hectáreas (sin contar los pastos permancns ni los cultivos permanentes)
    } Con el fin de fomentar el recambio generacional, el pago básico otorgado a los jóvenes agricultores menores de 40 años, nuevos agricultores o explotaciones fundadas en los últimos cinco años debe completarse con hasta un 25\% adicional del valor de los derechos durante los primeros cinco años de instalación. Las disposiciones relativas a la condicionalidad se confirman y simplifican. Para resolver el problema de los denominados «agricultores de sofá» y para eliminar las lagunas jurídicas que permitieron a un número limitado de empresas reclamar pagos directos, a pesar de que su actividad económica principal no era agrícola, la reforma impuso una lista negativa de actividades excluidas de recibir pagos directos, ampliable por los Estados miembros.

    o  Componentes voluntarios: Los Estados miembros pudieron utilizar hasta el 30\% de la dotación nacional y redistribuirla a los agricultores por sus 30 primeras hectáreas ( o hasta el tamaño medio de las explotaciones, si es mayor de 30 hectáreas) o aplicar un pago máximo  por hectáreá. Los Estados miembros pudieron conceder a los pequeños agricultores un régimen simplificado de un pago anual, independientemente del tamaño de la explotación, con requisitos de condicionalidad menos estrictos y exentos de la ecologización. Los Estados miembros o las regiones pueden destinar a zonas condicionadas por limitaciones naturales o zonas desfavorecidas un pago adicional de hasta el 5% de la  dotación nacional. A fin de solucionar posibles efectos perjudiciales de la convergencia interna para sectores o regiones especialmente sensibles, los Estados miembros podrán establecer pagos vinculados a productos específicos.


\subsection{Nueva PAC 2023-2027}
Tras la reforma de 201 3 el contexto económico e institucional se vio notablemente alterado. Si bien el crecimiento volvió, este siguió siendo muy bajo, antes de pasar a ser negativo como consecuencia de la COYID-19. La incertidumbre en los mercados se agravó tras el Brexit y la crisis en la tradicional asociación entre la UE y los EEUU a raíz de la elección del presidente Trump. El multilateralismo comercial inició un estancamiento y surgieron nuevos desafios en relación con el cambio climático, al mismo tiempo que entró en vigor el Acuerdo de París (COP2 l ). Las innovaciones técnicas, en particular la revolución digital, entrañaron importantes consecuencias para la producción, transformación y distribución de alimentos. En este contexto, los trabajos para adaptar la PAC a la nueva realidad institucional, económica y presupuestaria comenzaron en 20·1 6.

Ante la creciente preocupación sobre la evolución mundial de la cadena de suministro de alimentos y su impacto sobre la posición del productor en la misma, la Comisión creó en 20 16 un Grupo Operativo encargado de reflexionar acerca del futuro del Primer pilar con vistas a mejorar la posición de los agricultores en la cadena alimentaria. El Grupo fonnuló una serie de recomendaciones de actuación, entre las que cabe mencionar la intensificación de la transparencia del mercado, la prohibición de las prácticas comerciales desleales y algunos cambios complementarios de las normas de competencia que permitieran a los agricultores cooperar entre sí en el seno de las OP, que fueron adoptadas posteriormente por la Comisión. 

La Conferencia Europea sobre Desarrollo Rural Cork 2.0 de 20 16 adoptó la Declaración de Cork 2.0, "Una vida mejor en el medio rural", en la que se incluyeron diez orientaciones políticas para poder alcanzar una política de desarrollo rural innovadora, coherente e integradora que contribuyese a hacer frente a los retos económicos, sociales y medioambientales del medio rural. A partir de estas diez orientaciones y el resultado de la consulta pública sobre el futuro de la PAC de 2017, la Comisión presentó en 20 17 su Comunicación titulada «El futuro de los alimentos y de la agricultura». La Comunicación hace hincapié en el desarrollo sostenible, la preservación de los recursos naturales y la necesidad de garantizar el relevo generacional, así como en la gobemanza agrícola y anticipa un giro radical en el modelo de aplicación de la PAC.

Con motivo de la revisión intermedia del marco financiero plurianual 20 1 4-2020, la Comisión presentó en 2016, una propuesta legislativa Ómnibus que afectaba a numerosas políticas europeas, entre ellas la PAC. En un principio, el objetivo era realizar ajustes para simplificar los instrumentos existentes. Sin embargo, el Reglamento Ómnibus acabó convirtiéndose en una auténtica minirreforma de la PAC.\footnote{%
    Reglamento (UE) 2017/2393 del Parlamento Europeo del Consejo de 13 de diciembre de 2017.
} Las mejoras aprobadas se centraron en el ámbito de actuación de las OP, el refuerzo de los seguros agrarios y de los instrumentos de estabilización de los ingresos, las nonnas relativas a los pagos de ecologización y a los jóvenes agricultores, y la definición de «agricultor activo». 

Finalmente, la Comisión presentó sus propuestas legislativas para modernizar y simplificar la PAC en 2018. Los aspectos principales de las propuestas fueron los siguientes:
• pagos directos e intervenciones de desarrollo rural más selectivos y sujetos ambos a una planificación estratégica.
• una nueva arquitectura «verde» basada en condiciones medioambientales obligatorias y medidas voluntarias adicionales en el marco de ambos pilares.
• un planteamiento basado en los resultados (et «nuevo modelo de aplicación»), según el cual los Estados miembros tendrán que notificar sus logros cada año.

El retraso por los prolongados debates en el Parlamento Europeo y en el Consejo sobre las propuestas de la Comisión y el marco financiero plurianual, puso de manifiesto que no se podría aplicar a partir del 1 de enero de 2021 la nueva PAC. En consecuencia, la Comisión propuso en 20 J 9 una prórroga del marco jurídico existente y a finales de 2020 se adoptó el Reglamento Transitorio de la PAC  para garantizar la continuidad de la ayuda jurídica y financiera hasta finales de 2022, cuando entrase en vigor la nueva PAC.\footnote{%
    Reglamento (UE) 2020/2220 del Parlamento Europeo y del Consejo de 23 de diciembre de 2020.
}

La nueva PAC se adoptó formalmente en diciembre de 2021 con la publicación de los siguientes reglamentos:
• Reglamento de Planes estratégicos de la PAC: Reglamento (U E ) 202I /2 I 1 5 del Parlamento Europeo y del Consejo de 2 de diciembre de 202 1 .
• Reglamento de financiación, gestión y seguimiento de la PAC: Reglamento ( U E ) 2021/2 1 1 6 del Parlamento Europeo y del Consejo ele 2 de diciembre de 202 1 .
• Reglamento sobre la organización común de los productos agrarios (OCM): Reglamento ( U E ) 2021/2 1 1 7 del Parlamento Europeo y del Consejo de 2 de diciembre de 2021

Si bien se mantienen sus elementos esenciales. la PAC pasa de ser una política basada en la descripción de los requisitos que deben cumplir los beneficiarios de las ayudas, a una política con un fuerte énfasis en la consecución de resultados concretos y el desempeño. La nueva PAC será, además, una herramienta clave para alcanzar los objetivos del Pacto Verde Europeo. El Pacto Verde Europeo es la hoja de ruta de la Comisión para lograr que la economía de la UE sea sostenible. Dentro de las acciones clave del Pacto Verde Europeo, se encuentran dos eón una gran influencia sobre el sector agroalimentario, la Estrategia "De la Granja a la Mesa" y la Estrategia de la UE sobre Biodiversidad para 2030. La Estrategia "De la Granja a la Mesa" para un sistema alimentario justo, saludable y respetuoso con 71 medio ambiente pretende armonizar el sistema alimentario con las necesidades del planeta para obtener alimentos sanos, equilibrados y respetuosos con el clima y el medio ambiente. La estrategia de la UE sobre Biodiversidad para 2030 parte de la premisa de que proteger y recuperar la biodiversidad y el buen funcionamiento de los ecosistemas es fundamental para reforzar la resiliencia y prevenir la aparición y propagación de enfermedades en el futuro. Además de que invertir en la protección y recuperación de la naturaleza será también fundamental para la recuperación económica de Europa tras la crisis de la COVID- 1 9. 

La Estrategia "De la Granja a la Mesa" aborda de manera integral los desafios a los que se enfrentan los sistemas alimentarios sostenibles a partir de un Plan de Acción con 27 medidas aplicables desde 2020 hasta 2024. Estas medidas se pueden agrupar en 6 grupos: ✓ Medidas para garantizar la seguridad alimentaria. ✓ Medidas para estimular las prácticas sostenibles de elaboración de alimentos, venta al por mayor y al por menor, hostelería y servicios alimentarios. ✓ Medidas para promover el consumo sostenible de ali¡nentos y facilitar el cambio a dietas saludables y sostenibles. ✓ Medidas para reducir la pérdida y el desperdicio de alimentos. •✓ Medidas de lucha contra el fraude alimentario a lo largo de la cadena de suministro de alimentos. ✓ Medidas para garantizar la producción sostenible de alimentos. Los compromisos inclüidos dentro de la estrategia de la UE sobre Biodiversidad para 2030 relacionados directamente con la nueva PAC son: ✓ Invertir la tendencia a la disminución de los polinizadores. ✓ Al menos el 10% de la superficie agrícola contendrá elementos de paisaje de alta diversidad como los márgenes multifuncionales, muros, terrazas, charcas, etc. Ambas estrategias comparten los siguientes objetivos comunes: ✓ Reducir para 2030 en un 50% el uso y riesgo de pesticidas de síntesis químicos y en un 50% el uso de los pesticidas de alto riesgo. ✓ Reducción del 50% del exceso de nutrientes (especialmente fósforo y nitrógeno) y del 20% del uso de fertilizantes. ✓ Reducir en u�� 50% las ventas de antimicrobianos para los animales de granja y en acuicultura. ✓ Al menos el 25% de la superficie agraria europea deberá ser de agricultura ecológica.

\subsubsection{Objetivos}
La nueva PAC se centra en la consecución de la sostenibilidad social, medioambiental y económica de la agricultura y las zonas rurales. Los tres pilares de la sostenibilidad son: • fomentar un sector agrícola inteligente, competitivo, resiliente y diversificado que garantice la seguridad alimentaria a largo plazo. • apoyar y reforzar la protección del medio ambiente y contribuir a alcanzar los objetivos medioambientales y climáticos de la UE, entre ellos los compromisos contraídos en virtud del Acuerdo de París. • fortalecer el tejido socioeconómico de las zonas rurales. Estos objetivos generales se desglosan a su vez en nueve objetivos específicos y un objetivo transversal común.

Objetivos específicos
1. Apoyar unos ingresos agrícolas viables y la resiliencia del sector con el fin de mejorar la seguridad alimentaria a largo plazo y la diversidad agrícola, así como garantizar la sostenibilidad económica de la producción.
2. Mejorar la orientación al mercado y aumentar Ja competitividad agrícola con una mayor atención a la investigación, la tecnología y la digitalización.
3. Reequilibrar el poder de la cadena alimentaria y mejorar la posición de los agricultores en la misma.
4. Contribuir a la mitigación y adaptación al cambio climático, así como la promoción de la energía sostenible. 5. Fomentar el desarrollo sostenible y la gestión eficiente de los recursos naturales.
6. Contribuir a detener y revertir la pérdida de biodiversidad, mejorar los servicios ecosistémicos y preservar hábitats y paisajes.
7 . Atraer y sostener a los agricultores jóvenes y nuevos agricultores y facilitar el desarrollo empresarial sostenible en las zonas rurales.
8. Promover el empleo, el crecimiento, la igualdad de género, la inclusión social y el desarrollo local en las zonas rurales, así como la bioeconomía circular y la silvicultura sostenible.
9. Mejorar la respuesta a las demandas sociales de alimentos de alta calidad, seguros y producidos de manera sostenible, para reducir el desperdicio de alimentos, así como para mejorar el bienestar animal.

Objetivo transversal
1 0. Modernizar la agricultura y las zonas rurales mediante el fomento y el intercambio de conocimientos, la innovación y la digitalización.


\subsubsection{Implementación: Nueva Gobernanza}

✓ Planificación estratégica para cada Estado miembro
Cada Estado miembro diseñará un Plan Estratégico en el que indicará las intervenciones con las que alcanzar los objetivos de la PAC y el Pacto Ver��e Europeo. El contenido de los planes debe abarcar un análisis DAFO de las necesidades específicas del sector agrario y el medio rural, la definición de las medidas para dar respuesta a las necesidades existentes, un plan de metas para sus propios objetivos específicos y un plan financiero. El Plan Estratégico debe combinar ayudas directas, medidas de desarrollo rural y medidas sectoriales. De esta fonna, cada país elije las intervenciones específicas que considere más eficaces en función de sus objetivos concretos, basándose en una evaluación clara de sus propias necesidades. Cada Plan Estratégico requiere la aprobación previa de la Comisión para garantizar su coherencia con los objetivos de toda la UE, en concreto con la legislación y los compromisos medioambientales y climáticos de la U E y que no distorsionan el mercado único ni producen cargas excesivas para los beneficiarios o las administraciones.

✓ Enfoque en el rendimiento y los resultados El Reglamento de los planes estratégicos de la PAC fija un Marco de Rendimiento, Seguimiento y Evaluación (MRSE) para el período 2023-2027. Este marco abarca ambos pilares Y. supone un cambio de orientación estratégica que pasa del cumplimiento de las normas al rendimiento y los resultados. Este nuevo modelo utiliza un conjunto de indicadores comunes de rendimiento: o de . realización, que se utilizarán para el seguimiento de la aplicación de la PAC. o de resultados, que se utilizarán para supervisar los avances de los. países hacia los objetivos fijados. o de contexto y de impacto, que se utilizarán para evaluar el rendimiento global de las políticas respecto de los objetivos de la PAC.

Los indicadores se supervisarán mediante infonnes anuales de rendimiento y una revisión semestral de los resultados de los planes estratégicos para evaluar los progresos de los países en la consecución de los objetivos.

\subsubsection{Medidas principales}
Una PAC más verde: Los dos pilares se combinan y orientan hacia los mismos objetivos ambientales y climáticos, mostrando una mayor ambición ecológica.
• Los planes estratégicos deberán estar en consonancia con la legislación medioambiental y climática. Cada país estará obligado a mostrar una mayor ambición en materia medioambiental y deberá actualizar el plan estratégico cuando se modifique la legislación en materia de clima y medio ambiente. 

• Las nuevas nonnas de condicionalidad incluirán un mayor nivel de ambición. Las reglas que se espera que cumplan los agricultores incluyen: o Los requisitos legales de gestión (RLG), aplicables a todos los agricultores, que engloban nonnas de la U E sobre salud pública, animal y vegetal, bienestar de los animales y el medio ambiente. o Buenas condiciones agrícolas y medioambientales (BCAM), aplicables únicamente a los agricultores que reciben apoyo de la PAC. Estos estándares están diseñados para prevenir la erosión del suelo, mantener su materia orgánica y estructura, conservar pastizales pennanentes, proteger la biodiversidad y garantizar la conservación de las características del paisaje, prot��ger y gestionar el agua mediante el establecimiento de franjas de amortiguamiento a lo largo de los cursos de agua, la autorización de agua para riego y la protección de las aguas subterráneas contra la contaminación.

• Al menos el 25 % del presupuesto para pagos directos se asignará a los-.esquemas ecológicos que los países deberán incluir de manera obligatoria en sus planes, pero que serán de voluntario cumplimiento para los agricultores. Las primas pueden establecerse como pagos de '·compensación" por los costos adicionales y la pérdida de ingresos derivados de las prácticas en cuestión o como pagos que van más allá de la compensación, siempre que se respeten las reglas pertinentes de la Caja Verde. Durante un "periodo de aprendizaje" de dos años, un país puede destinar menos del 25% del presupuesto en caso de que la superficie potencial que se acoja a una práctica sea menor de lo previsto, siempre que compense la mayor parte del déficit a finales de 2027.

• El segundo pilar aumentará la parte de sus recursos destinados a intervenciones relacionadas con el clima y el medio ambiente: o al menos el 35 % de los fondos (frente al 30% actual) se destinará a compromisos de gestión agroambiental, Natura 2b00 y pagos de la Directiva Marco del Agua, inversiones ambientales y climáticas, y bienestar animal. o el 50 % de los pagos irá destinado a zonas con limitaciones naturales.

• El 40 % del presupuesto tendrá que destinarse a !}Ctuaciones pertinentes para el clima y apoyar el compromiso general de dedicar el 10 % del presupuesto de la UE a los objetivos de biodiversidad antes de que finalice el período del marco financiero plurianual.

✓ U na PAC más ju sta: se dirige el apoyo a quienes más lo necesitan mediante una distribución más equitativa de los pagos:

• Los países tendrán que dedicar al menos el I O % de sus asignaciones de pagos directos a un mecanismo redistributivo de ayuda a la renta para abordar mejor las necesidades de ingresos de las explotaciones pequeñas y medianas. Los países podrán establecer excepciones a esta obligación si las necesidades redistributivas se abordan adecuadamente a través de otros medios, como, por ejemplo, la reducción de los pagos, - la convergencia interna o la territorialización de las ayudas a la renta básica. La reducción y limitación de los pagos directos para las explotaciones que reciben grandes cantidades será opcional. Se permitirá a los Estados miembros reducir hasta un 85 % los pagos directos por agricultor o ganadero que superen los 60 000 euros al año y limitarlos a 100 000 euros al año. El apoyo a las pequeñas explotaciones se reforzará con la posibilidad de sustituir los diferentes pagos directos por una única ayuda. Los países podrán seguir concediendo una parte limitada de su dotación de'pagos directos para apoyar a algunos sectores o tipos de agricultura. Las ayudas acopladas tendrán
como objetivo hacer frente a las dificultades específicas de esos sectores y tipos de agricultura
mejorando su calidad, sostenibilidad o competitividad.

• La nueva normativa contiene una definición obligatoria, pero flexible de agricultor activo, que se basará en criterios objetivos y no discriminatorios, como la prueba de ingresos, el objeto social y la inclusión en un registro. Los países deben garantizar que los agricultores pluriactivos y a tiempo parcial no queden excluidos de las ayudas y podrán determinar que los agricultores que reciben pagos directos de hasta 5.000€ se consideren agricultores activos.

• Por primera vez los pagos estarán vinculados al respeto de determinadas disposiciones de la legislación laboral de la UE relativas a las condiciones de empleo, la seguridad y la salud en la explotación agrícola.

• Los nuevos niveles de ayuda a la renta convergerán más, tanto dentro de los distintos países (convergencia interna) como entre los países (convergencia externa). Por la convergencia interna, los países que todavía realizan pagos directos basados en referencias históricas deberán seguir reduciendo las diferencias entre los pagos por hectárea. Este proceso debe garantizar que de aquí a 2027 todos los derechos de pago alcancen un valor mínimo del 85 % del importe medio nacional. Por el proceso de convergencia externa, los países cuyo nivel de ayuda directa por hectárea sea inferior al 90 % del pago medio por hectárea de la UE verán satisfecha la mitad de esta diferencia para 2027.

• Los países tendrán que distribuir al menos el 3% de su dotación de pagos directos para apoyar a los jóvenes agricultores que establezcan una explotación agrícola. Esta ayuda puede otorgarse como apoyo a la renta, apoyo a la inversión o ayudá inicial para jóvenes agricultores.

• La igualdad de género y el aumento de la participación de las mujeres en la agricultura forman parte, por primera vez, de los objetivos de la PAC.

✓ Una PAC más com petit iva: se reforzará la posición de los agricultores en la cadena de suministro e impulsará la competitividad del sector agroalimentario: 

• Se fortalece la posición negociadora de los agricultores a través de intervenciones sectoriales: o ampliación del apoyo a las OP. siguiendo el modelo de frutas y hortalizas, a todos los sectores agrícolas. excepto el vitivinícola y la apicultura. o todos los programas de apoyo a sectores específicos actualmente cubiertos por la OCM formarán parte de los planes estratégicos. o se amplían algunas excepciones de la ley de competencia. lo que permite a los agricultores trabajar juntos y/o con otras partes interesadas en la cadena alimentaria con el fin de garantizar estándares de sostenibil idad más altos. o se extenderá el actual instrumento de regulación de la oferta para las DOP y las IGP a todos los sectores. hasta ahora limitados al jamón y el queso.

• se mantiene la orientación general al mercado de las reformas anteriores, reforzada con el apoyo de los nuevos planes estratégicos y las normas de la OCM. Se crea una nueva reserva financiera de al menos 450 millones de euros cada año. 

• se han acordado normas específicas para mejorar el apoyo al sector vitivinícola. o se prorroga el régimen de autorización de plantaciones de vid hasta 2045. o se autoriza utilizar variedades de vid híbridas de "Yitis vinifera" y otras especies de "Yitis" para la producción de vinos de calidad amparados por DOP. o se autoriza comercializar vinos desalcoholizados y proteger los productos vitivinícolas parcialmente desalcoholizados como DOP e IGP. o se implementan nuevas reglas sobre el etiquetado de ingredientes y nutrición para vinos y productos vitivinícolas aromatizados para mejorar la información del consumidor, exigiendo a los productores de vino que indiquen el valor energético de sus productos, y asegurando que se proporcione información completa sobre nutrición e ingredientes en la etiqueta o por medios electrónicos.

• se establecen normas adicionales relativas a las IG en consonancia con algunas de las medidas adoptadas para reforzar la posición de los agricultores y las nuevas normas para el sector vitivinícola: o los criterios de procesamiento sostenible pueden incluirse en la especificación del producto de forma vol untaría. o las DOP y las IGP estarán expresamente protegidas en Internet y contra las mercancías en tránsito dentro de la UE. La protección de las DOP e IGP también se extenderá al uso de los nombres de los ingredientes. o los procedimientos se simplificarán, agilizarán y se harán coherentes entre sectores.

\subsubsection{Financiación}
El marco funanciero plurianual para el período 202 1 -27 prevé una asignación total para la PAC de 386.600 millones de euros. El porcentaje que los gastos agrícolas representan en el presupuesto de la UE mantiene su tendencia descendente y para el período 202 1-2027 representan el 3 1 % frente al 37,8% del período anterior. Durante los dos primeros años del marco financier plurianual se aplicaron las disposiciones contenidas en el Reglamento Transitorio de la PAC entre las que cabe destacar el desembolso del 30% de los 8.070 millones de euros del Instrumento Europeo de Recuperación en 202 1 y el 70 % en 2022, el desembolso de un tercio del presupuesto total (37%) para medidas ecológicas y de bienestar animal y más de la mitad del presupuesto total (55%) para medidas sociales y de transfonnación digital y la prórroga de la aplicación de la ayuda temporal excepcional a los agricultores y las PYME afectadas por la COY! D-19 por seis meses.

El FEAGA cuenta con una dotación de 291 . 100 millones de euros, de los cuales hasta 270.000 millones de euros se destinarán a planes de apoyo a la renta, y el resto se dedicará a apoyar los mercados agrícolas. Para el FEADER, la dotación total asciende a 95.500 millones de euros. Esto incluye 8. 100 millones de euros del plan Next Generation EU para ayudar a abordar los desafíos planteados por la pandemia de COVID- 19 y a las zonas rurales a realizar los cambios estructurales necesarios para alcanzar los objetivos del Pacto Verde Europeo y la transición digital. Para que los países puedan adaptar mejor la utilización de los fondos a las prioridades de sus sectores agrícolas, tendrán derecho a transferir hasta el 25 % de sus asignaciones entre la ayuda a la renta y el desarrollo rural.

La nueva PAC empezará a aplicarse el I de enero de 2023, una vez finalice la vigencia del Reglamento Transitorio de la PAC y comiencen a aplicarse los Planes Estratégicos. A fines de 2023 la Comisión presentará un informe para evaluar el esfuerzo conjunto de todos los Planes Estratégicos de la PAC, prestando especial atención a la ambición colectiva para alcanzar los objetivos del Pacto Verde Europeo. A partir de 2024, cada país presentará un informe anual de rendimiento y celebrará una reunión de revisión anual con la Comisión. La Comisión llevará a cabo una primera revisión del rendimiento de cada Plan Estratégico de la PAC y solicitará, si fuese necesario, acciones de seguimiento específicas a los países de la UE. En 2026, una evaluación intermedia evaluará el desempeño de la nueva PAC. La Comisión llevará a cabo en 2027 una segunda revisión del rendimiento de cada Plan Estratégico de la PAC.












\clearpage
\section{Política Pesquera Común}
\puntosclave{
    La PPC obedece al significativo impacto a nivel social, económico y medioabiental del sector pesquero.
    
    Al igual que la PAC, ha ido evolucionando para hacer frente a una realidad cambiante. 
    
    Desde 2013 está orientada a garantizar la sostenibilidad medioambi ental, económica y social a largo plazo del sector pesquero así como su competitividad y transparencia. 
    
    Dado que el 13\% de todas las capturas procede de aguas internacionales y el 8\% de las ZEE, su dimensión exterior es muy relevante.
}

La PPC, de importancia menor en comparación con la PAC, adquirió relevancia de manera progresiva debido a la importancia económica del sector de la pesca en la economía de algunos países de la UE y las limitaciones de acceso a los recursos pesqueros derivados de la normativa internacional que motivaron la necesidad de articular una política que garantice un adecuado acceso a los recursos marinos y una eficiente gestión y conservación de los mismos.



\subsection{¿Qué es la Política Pesquera Común? Naturaleza y objetivos}
La UE dispone de un sector pesquero reducido económicamente, pero que disfruta de una alta prioridad política. Aunque su contribución al PNB es pequeña, el impacto de la pesca y de la acuicultura es altamente significativo como fuente de empleo en regiones que a menudo no ofrecen otras alternativas económicas. Adicionalmente, se trata de un sector primario que contribuye a asegurar el consumo alimentario de los ci udadanos y aunque las poblaciones de peces, por lo general tienen una capacidad reproductiva alta, esta no es ilimitada. Si la pesca no se controla, las poblaciones pueden extingui rse o dejar de ser económicamente viable.

La política pesquera común (PPC) fue instaurada por el Tratado de Roma. Inicialmente integrada en la PAC y con una relevancia residual, con el tiempo se configuró como una política independiente. El fundamento jurídico de la PPC se encuentra recogido en los artículos 38 a 43 del TFUE y al igual que ocurre con la PAC, se trata de una competencia compartida.

Como política integrada en la PAC, la PPC compartía sus objetivos. pero de manera específica, los objetivos de la PPC en su origen eran preservar las poblaciones de peces, proteger el medio ambiente marino, garantizar la viabilidad económica de las flotas europeas y proporcionar a losconsumidores alimentos de calidad.

\subsection{Desarrollo histórico}

El Tratado de Roma estableció el libre acceso a las aguas comunitarias de la flota comunitaria, de tal forma que los buques pesqueros tuviesen igualdad de acceso a las aguas y a los recursos marinos comunitarios. Los seis miembros fundadores contaban con un sector pesquero déb il y la actividad pesquera se desarrollaba básicamente entorno a sus puertos base, por lo que este principio no se aplicó y tampoco se desarrolló una política específica en ese momento. Sin embargo, a partir de 1970 la política pesquera fue adquiriendo de manera progresiva. una identidad propia que culminó en 1983 con la implementación de una PPC propiamente dicha. Los factores que motivaron este desarrollo fueron el reconocimiento a nivel mundial de las "zonas económicas exclusivas" (ZEE) y la adhesión a la Comunidad de países que disponían de flotas pesqueras importantes como Reino Unido, Irlanda y Dinamarca lo que dotó de una nueva dimensión al sector pesquero común.\footnote{%
    Espacio marítimo que se extiende a una distancia de 200 millas desde la costa en el que los Estados ribereilos ejercen una soberanía fu ncional y jurisdiccional en lo relativo a la exploración y explotación de los recursos naturales. de las aguas. del lecho y el subsuelo marino. Se establecieron en la 111 Conferencia de las Naciones Unidas sobre el Derecho del Mar de 1982
} Estos factores obligaron a la Comunidad a dar respuesta a problemas específicos de la pesca como el acceso a los recursos comunes, la conservación de las poblaciones, las medidas estructurales para la flota pesquera y las relaciones internacionales.

\subsubsection{Configuración inicial: 1983-2002}
En esta primera etapa se articularon cuatro grandes políticas:

1. Política de recursos, relativa al acceso a los recursos marinos comunitarios y la gestión y conservación de los mismos. El acceso se basó en el principio de libre acceso. Este principio es de aplicación en las aguas comunitarias comprendidas entre las primeras 12 millas náuticas y las 200 millas y fundamentalmente, en las agu��s comunitarias atlánticas, debido a que la estrechez de la plataforma continental en el Mediterráneo hace que la gran mayoría de las zonas de pesca se encuentren dentro de las primeras 12 millas.\footnote{%
    Las primeras 12 millas náuticas desde la costa son aguas territoriales de los Estados, sobre las que tienen competencia exclusiva y pueden restringir el acceso a la flota pesquera de los otros Estados miembros. Esta excepción al principio de libre acceso se basa en la necesidad de preservar estas zonas. más sensibles, mediante la limitación del esfuerzo pesquero y la protección de las actividades pesqueras tradicionales de las que depende el desarrollo social y económico de detenninadas comunidades costeras.
} Respecto a la gestión y conservación, se fijó como objetivo principal garantizar la viabilidad del sector a largo plazo mediante una explotación sostenible de los recursos. Se configuró un sistema basado en el cálculo anual de los Totales admisibles de capturas (TACs) para cada especie protegida y para cada zona marítima del Atlántico. Los TACs eran posteriormente repartidos entre los Estados miembros en forma de cuotas nacionales de acuerdo al principio de estabilidad relativa (PER), basado especialmente en los niveles históricos de capturas, quedando determinadas las posibilidades máximas de pesca de cada país. Este sistema implantó una distribución automática, permanente y fija, de las oportunidades de pesca independientemente de la competitividad de las flotas de cada país. El esfuerzo de pesca también se limitó anualmente fijando un número máximo de días de pesca por zona y tipo de embarcación y prohibiendo la pesca en determinados períodos y/o zonas para permitir la recuperación de las especies.\footnote{%
    El esfuerzo de pesca se define como la capacidad de pesca (medida en toneladas de registro bruto o potencia instalada de los barcos) multiplicada por la actividad pesquera. expresada en días pasados en el mar. La gestión del esfuerzo pesquero consiste en combinar las limitaciones de la capacidad de la flota y el tiempo que esta pasa en el mar.
}

2. Política de estructuras. Su aspecto más relevante fue la financiación de la reducción de la flota comunitaria para eliminar el exceso de capacidad pesquera dada la disminución de las posibilidades de pesca. En 1992 se decidió incorporar la política estructural de la pesca a los Fondos Estructurales y crear el Instrumento Financiero de Orientación de la Pesca (IFOP). Este fondo, además de financiar las medidas estructurales, tenía como objetivo reducir la dependencia económica de aquellas regiones que viviesen exclusivamente de la pesca, fomentado actividades alternativas.

3. Política de mercados. Articulada a través de la correspondiente organización común de los mercados (OCM) para la que se fijaron cuatro objetivos:
• Estimular el asociacionismo a través de las Organizaciones de productores (OP, asociaciones de empresas pesqueras para regular los desembarcos y los precios) y las Organizaciones Interprofesionales de Productores (OIP, personas jurídicas que reúnen la representación de actividades vinculadas a la producción, a la transformación o al comercio con representación en una o más regiones europeas)
• Implantar un sistema de sostenimiento de precios, que estableciese precios mínimos por debajo de los cuales no se pudiesen vender los productos pesqueros. Estos precios de retirada se fijaban en función de la media de los precios registrados durante los tres aftos precedentes en los puertos representativos y eran precios de referencia, no de garantía, puesto que los organismos de intervención nunca compraban directamente al productor, este podía elegir entre la retirada y vender a bajo precio. También se previeron ayudas a las OP para retirar productos o almacenarlos.
• Elaborar una normativa común de productos frescos, que desarrolló un proceso de trazabilidad, categorías de empacado y calidad de los productos comunitarios e importados.
• Desarrollar un régimen de comercio con terceros países que garantizase la estabilidad de precios y unas condiciones de competencia leales.

4. Política exterior que se articuló a través de la firma de Acuerdos bilaterales y multilaterales de pesca. ·se alcanzaron diferentes acuerdos bilaterales en función de la capacidad de explotación de los recursos del tercer país.\footnote{%
    La pérdida de acceso a zonas tradicionales de pesca como consecuencia de la ampliación de las ZEE, el exceso de capacidad de la flota y el desarrollo de nuevas tecnologías de 'pesca industrial como barcos-factoría y grandes congeladores que pcm1iten pescar lejos de los puertos de origen. motivaron la necesidad de alcanzar acuerdos bilaterales.
} Acuerdos basados en intercambios recíprocos de derechos de pesca y Acuerdos de compensación por acceso al recurso, firmados normalmente con países que no tienen capacidad para explotar por sí mismos los recursos marinos queposeen y reciben una compensación financiera por permitir pescar en sus aguas. Los Acuerdos multilaterales se concretan en la participación en Organizaciones regionales de pesca (ORP) y en Convenios internacionales. Las ORP se crean para gestionar los recursos en alta mar o los recursos compartidos por diversos países para reforzar la coope��ación regional. Asumidos por la UE, vinculan a todos los Estados miembros, pero jurídicamente sólo comprometen a los Estados firmantes. Los Convenios internacionales ofrecen un marco jurídico a los mares  y océanos, promueven su uso pacífico, la utilización equitativa y eficaz de sus recursos, la conservación de sus recursos vivos y la protección y conservación del medio ambiente marino.


\subsubsection{Reforma 2002: }
a) Objetivos perseguidos
• Refonnular algunos aspectos de la política de gestión y conservación. El PER y el sistema de T ACs desincentivaban la modernización y la competencia, haciendo del sector ·esquero comunitario un sector cada vez menos competitivo y alejado de los parámetros mundiales. 

• Llevar a cabo un control más eficaz de la actividad para reprimir prácticas fraudulentas.\footnote{%
    Pesca en detenninadas áreas sin la autorización necesaria. o por encima de las cuotas asignadas. no respetar los paros biológicos. uso de redes y técnicas inadecuadas
}

• Superar la contradicción de que las cuotas se fijasen a nivel comunitario y el control se delegara a los Estados miembros.

• Hacer frente al exceso de capacidad de la flota comunitaria y la sobreexplotación
continuada de algunas especies.
b) Principales medidas

• Respecto a la política de recursos, se mantuvieron el PER y el sistema de cuotas, pero se permitió a los Estados distribuir las cuotas de la forma que considerasen más conveniente. Se introdujo un enfoque plurianual de la política de conservación a través de los planes de recuperación para las poblaciones de peces en peligro y planes de gestión de stocks para las demás poblaciones. En estos planes se fijaban objetivos de capturas plurianuales con la finalidad de evitar las modificaciones repentinas de los límites de capturas de un año para otro y aportar unas condiciones más estables de funcionamiento. (Ver Anexo 4: \textit{Granja en la mesa}, antes)

• Respecto a la Política de estructuras se estableció un nuevo sistema basado en unos niveles de referencia de capacidad de pesca expresados en toneladas de registro bruto o KW de potencia instalada. Se introdujeron medidas socioeconórnicas como ayudas paralos pescadores y propietarios de buques que interrumpiesen temporalmente su actividad por circunstancias imprevistas o ayudas a los pescadores para emprender actividades alternativas. En 2004 se creó el Fondo Europeo de Pesca (FEP) que sustituyó al IFOP.

• Se definió con mayor claridad la división de competencias entre los Estados miembros y la Comisión. Los Estados miembros vieron reforzada su responsabilidad en el control,  inspec;ción y cumplimiento de la política y se impulsó la cooperación y coordinación de los Estados en relación a las tareas de inspección. Se facultó a la Comisión a realizar auditorías, verificaciones e inspecciones para garantizar la aplicación de la PPC.

• En Política exterior, los Acuerdos bilaterales con contrapartida financiera se sustituyeron por los Acuerdos de asociación pesqueros (AAP). La flota accedía al excedente de pesca de ZEE pertenecientes, el) su mayor parte, a países de ACP y en compensación los países recibían un canon abonado por los armadores y un.importe a tanto alzado por parte de la UE. La contribución financiera de la UE se justificaba por el interés mutuo en invertir en una política sostenible y no constituía un mero pago por derechos de acceso.


\subsection{PPC en la actualidad: Reforma de 2013}
La reforma de 2002 no detuvo el deterioro de algunas poblaciones y dejó al descubierto algunos problemas que hasta entonces no se habían detectado, corno el de los descartes.\footnote{%
    El descarte es la práctica de devolver al mar las capturas no deseadas, vivas o no. por no alcanzar la talla, porque el pescador no dispone de cuota o por determinadas nonnas de composición de las capturas.
} En 2013 se alcanzó un acuerdo para un nuevo régimen de pesca que entró en vigor en 20 14 y que se basa en cuatro pilares fundamentales:

\subsubsection{Política Pesquera Común}
Objetivo: (Reglamento (UE) nº 1 380/201 3 sobre la Política pesquera común.) garantizar que las actividades en los sectores de la pesca y la· acuicultura contribuyan a la sostenibilidad medioambiental, económica y social a largo plazo. En consonancia con el Pacto Verde Europeo y la Estrategia sobre la biodiversidad para 2030, laspesquerías de la UE se rigen por el principio de cautela para limitar el impacto de las actividades pesqueras sobre el ecosistema marino.

✓ Principales medidas
• Gestión plurianual ecosistémica que refuerce la importancia de los planes plurianuales, pero prestándo más atención a los ecosistemas y sustituyéndo los planes para especies concretas por planes para varias especies y caladeros.
• Se introduce el principio de Rendimiento máximo sostenible (RMS), rendimiento de equilibrio teórico máximo que puede extraerse continuamente, en promedio. d e una población en las condiciones ambientales medias existentes sin que ello afecte significativamente al proceso de reproducción. Según este principio, el nivel de cuotas fijado será aquel que permita alcanzar el RMS.
• Se introdujo la p��ohibición de los descartes y la obligación de desembarcar todas las capturas y deducirse de las cuotas.
• Los Estados miembros elaborarán planes para reducir las capacidades siempre que se produzca un exceso de capacidad en cualquier segmento de flota.
• La pesca costera adquiere especial relevancia, dado su peso en el sector pesquero de la UE. La zona de exclusión de 1 2 millas náuticas establecida en favor de las flotas tradicionales seguirá vigente hasta 2022 y se recomendará a los Estados miembros que asignen una mayor parte de sus cuotas a ese sector, visto su escaso impacto medioambiental y elevada intensidad de mano de obra.
• Para aumentar el rendimiento y potenciar el crecimiento en las zonas costeras y rurales se elaborarán planes nacionales que supriman barreras administrativas y apliquen las normas medioambientales, sociales y económicas en el sector de la acuicultura.
• Se refuerza el aspecto científico, incrementándose la recogida de datos y la puesta en común de información sobre poblaciones, flotas e impacto de la actividad pesquera.
• Para acercar el procedimiento de toma de decisiones a los caladeros, el marco general se definirá a nivel comunitario y los Estados m·iembros desarrollarán las medidas de aplicación

\subsubsection{Organización Común de Mercados}
Reglamento (UE) nº 1 3 79/20 1 3 por el que se establece la nueva Organización común de mercados (OCM) en el sector de los productos de la pesca y de la acuicultura.

El objetivo de la reforma es reforzar la competitividad y mejorar la transparencia mediante la modernización y simplificación de la normativa existente. Las OP cobran mayor protagonismo por lo que se refiere a la gestión, seguimiento y control colectivos. Para contribuir a la sostenibilidad, las OP y AOP establecen planes de producción y comercialización que penniten la gestión colectiva de las actividades de los productores. También se establecieron nuevas normas de comercialización en cuestiones como el etiquetado, la calidad y la trazabilidad. Las normas comunes de comercialización establecen unas características uniformes para los  productos de la pesca que se venden en la UE, sea cual sea su origen. Se aplican de acuerdo a medidas de conservación y contribuyen a la transparencia del mercado y a la calidad de los productos.

\subsubsection{Política Exterior}
Se introdujeron los Acuerdos de colaboración en materia de pesca sostenible (ACPS) que pretenden reforzar la gobemanza y el desarrollo social y económico de los países socios en el sector pesquero. Estos acuerdos permiten el acceso a los recursos a cambio de una contribución financiera y de un apoyo técnico que deben ayudar a lograr un seguimiento, un control y una supervisión más e��caces. El acceso a las flotas se negocia para asegurar que las poblaciones se explotan de forma sostenible, teniendo en cuenta los principios de precaución y RMS y favoreciendo el acceso prioritario de la flota del tercer país. Actualmente hay 13 ACPS en vigor con los países socios en los océanos Atlántico, Índico y Pacífico. Más de 135 millones EUR del presupuesto de la UE se asignan a los ACPS, de los que, de media, un 46 % se destinan a apoyar el desarrollo sostenible del sector pesquero en los países socios, así como a aumentar su capacidad global de gobemanza en materia de pesca.

\subsubsection{Fondo Europeo Marítimo, de Pesca y Acuicultura (FEMPA)}
Reglamento (UE) nº 2021/1 139 del Parlamento Europeo y del Consejo de 7 de julio de 2021

El FEMPA es el nuevo fondo de las políticas marítima, pesquera y acuícola propuesto para el periodo 202 1-2027, en su��titución del anterior Fondo Europeo de Pesca (FEMP) y dotado con 6 140 millones de euros a precios corrientes. Se articula entorno a cuantro prioridades:
l . Fomentar la pesca sostenible y la recuperación y conservación de los recursos biológicos acuáticos.
2. Fomentar las actividades sostenibles en la acuicultura y la transformación y comercialización de productos de la pesca y la acuicultura.
3. Permitir una economía azul sostenible en ·J·as zonas costeras, insulares e interiores, y fomentar el desarrollo de las comunidades pesqueras y acuícolas.
4. Reforzar la gobernanza internacional de los océanos y permitir que los mares y océanos sean seguros, protegidos, limpios y estén gestionados de manera sostenible.







Los planes de recuperación, de carácter prioritario, establecían los objetivos de recuperación (tamafio de la pob lación y/o producción a largo plazo y/o tasa de mortalidad por causa de la pescay/o estabilidad de capturas) y el plazo previsible para su consecución. Para poder lograr el objetivo  de recuperación, se introdujo la posibilidad de adoptar medidas ya existentes (limitación de capturas o del esfuerzo de pesca, medidas de carácter técnico) con la novedad de poder adoptarlas de manera conjunta para una o varias especies y con carácter plurianual.

Los planes de gestión establecían una serie de medidas técnicas, en función de las características de las poblaciones y las pesquerías (especies objetivo, artes empleados, estados de las pob lacionesafectadas) y el impacto económico de las medidas aplicadas a las pesquerías de que se tratase. 

Ante evidencia de amenaza grave para la conservación de los recursos marinos y si fuese necesario tomar medidas inmediatas, se otorgó a la Comisión la posibilidad de implementar medidas de emergencia para un periodo máximo de seis meses, prorrogab les a otros seis más. Los Estados miembros podían aplicarlas respecto de sus propias aguas durante tres meses.

Se crearon los Consejos Consultivos Regionales (CCR) para aumentar el grado de participación de los pescadores y otros grupos de interés. Los CCR están integrados por pescadores, científicosy representantes del sector de la pesca y la acuicultura, grupos de defensa del medio ambiente y protección de los consumidores, y también pueden participar las autoridades nacionales y regionales de cualquier Estado miembro y la Comisión. Pueden ser consultados por la Comisión y presentar recomendaciones y sugerencias o informar de los problemas que suscite la aplicación de la PPC en su zona.

















\clearpage
\section{Conclusión} 


\subsection{Recapitulación}


\subsection{Extensiones y relación con otras partes del Temario}



\subsection{Opinión}

\subsubsection{Idea final (Salida o cierra)}

\clearpage
\pagenumbering{gobble}
\MyBibliography



\MyBibItem{DAWID y GATTI}{Chapter 2 - Agent-Based Macroeconomics}{2018}{https://doi.org/10.1016/bs.hescom.2018.02.006}


% ANEXOS
\clearpage
\anexo{\textbf{Preguntas Test}}
%\input{\TemaCode/questions.tex} 



\end{document}